\section{Gaussian estimates at the bottom edge scales}
\label{sec:edge-estimates}

The observation below reveals noisy information about leading spin
projections.  We control the variance of a bounded spin function left
unresolved by that observation by \(Cp^{-2}\) times its heat-bath energy,
with a small remainder, at every scale for which \(np^3\) exceeds a tight
threshold.  The exponent \(2\) is important: a fixed
relative covariance error at every scale would accumulate to a power loss.
We first extend the Gaussian cap estimates to these scales and then show
that the relative error is integrable with respect to \(\mathrm d\log p\).

\subsection{Cap observations and the estimates}

Throughout this section the disorder is Gaussian, with the independent
GOE diagonal and randomized eigenvector signs specified in Section~2.
Set \(r_n=n^{-1/3+10^{-4}}\), the terminal scale of the source estimates,
and write \(d=2I-W\).  Define, in physical coordinates,
\[
\begin{aligned}
 P_p&=U\diag(\ind_{\{d_a\le p^2\}})U^{\mathsf T},&
 \cH_p&=\operatorname{ran}P_p,\\
 B_p&=U\diag((p^2-d_a)_+)U^{\mathsf T},&
 W_p&=W-B_p,\\
 k_p&=np^3,&M_p&=\sqrt{np}.
\end{aligned}
\]
Thus \(p^2M_p^2=k_p\).  A field at scale \(p\) lies in \(\cH_p\), with
length unit \(p^2M_p\).  For such a field let
\[
 \phi_p(h)=\log\nu_n
 \exp\left(\tfrac12 X^{\mathsf T}W_pX+h^{\mathsf T}X\right),
 \qquad m_p(h)=\nabla\phi_p(h),\qquad
 K_p(h)=\left[P_p\nabla^2\phi_p(h)P_p\right]_{\cH_p}.
\]
Brackets with subscript \(\cH_p\) denote compression to that space.
The probability with this log partition function is \(\mu_{p,h}\).
Thus \(m_p\) and \(K_p\) are its mean and its active covariance.
Superscript \({\rm sp}\) denotes the same definitions with normalized
surface measure on \(\{X:\|X\|^2=n\}\) in place of \(\nu_n\).
In the tight-constant convention of Section~\ref{sec:background}, a tight
lower threshold for \(k_p\) means that, for each \(\chi>0\), a deterministic
threshold \(K_\chi\) can be chosen on an event of limiting lower
probability at least \(1-\chi\).

For \(X\sim\mu\), form the Gaussian observation
\begin{equation}\label{edge:observation}
 h_p=B_pX+\beta_p,\qquad
 \operatorname{Cov}(\beta_p,\beta_q)=B_{\min(p,q)}.
\end{equation}
One construction uses an independent scalar Brownian motion in spectral
coordinate \(a\), evaluated at \((p^2-d_a)_+\).  Bayes' formula gives
\(\Law(X\mid h_q,\ q\le p)=\Law(X\mid h_p)=\mu_{p,h_p}\) at fixed \(W\).
Write \(Q_\mu\) for the joint law of
the spin and all observations at fixed \(W\).  The observation filtration
excludes the spin.  The path \(h_p\) is continuous at \(p^2=d_a\); the
new coordinate then has zero precision.  The finitely many activation
points may be omitted in scale integrals.  For \(0\le f\le1\), put
\begin{equation}\label{edge:posterior-variance}
 F_p(h)=\mu_{p,h}[f],\qquad
 V_p(f)=Q_\mu[\Var(f\mid h_p)].
\end{equation}

For \(u>-p^2\) define
\[
 R_{p,u}=U\diag\left((\max\{d_a,p^2\}+u)^{-1}\right)U^{\mathsf T}.
\]
On the spectral events and scales used below, the spherical saddle
\(\alpha_p(h)>-p^2\) is the unique solution of
\begin{equation}\label{edge:saddle}
 n=\Tr R_{p,\alpha_p(h)}+\|R_{p,\alpha_p(h)}h\|^2.
\end{equation}
Existence, uniqueness, and bounds on this solution are proved below.
Since \(R_{p,u}\) is scalar on \(\cH_p\), the length of a field whose
saddle is \(u\) is
\begin{equation}\label{edge:saddle-length}
 \ell_p(u)=(p^2+u)\sqrt{n-\Tr R_{p,u}}
\end{equation}
whenever the expression under the square root is nonnegative.

We use one change of probability law, always conditional on a fixed
spectrum.  Under the ordinary law \(U\) is Haar, and \(Q_\mu\) is the
joint law of the spin and its physical observations at the resulting
\(W\).  With probability reference measures in the partition functions,
set \(L_\circ(U)=Z_{\rm cube}(W)/Z_{\rm sp}(W)\).
Haar averaging gives \(E_UL_\circ=1\).  Under \(\mathbf P\), first choose
\(U\) with Haar density \(L_\circ\), and then use \(Q_\mu\) given \(W\);
write \(\mathbf E\) for its expectation.  The spherical identity of
\cite[Lemma~2.5 and Equation~(2.8)]{FUA} gives
\begin{equation}\label{edge:haar-bias}
 \Law_{\mathbf P}(U^{\mathsf T}X)
   =\mu^{\rm sp}\text{ in spectral coordinates},\qquad
 E_U(L_\circ-1)^2=o(1)
\end{equation}
on its retained eigenvalue events.  Consequently the spectral observation
path \((U^{\mathsf T}h_p)_p\) under \(\mathbf P\) has the spherical
observation law, denoted by \(Q^{\rm sp}\).  For comparison, let
\(\mathsf R\) choose \(U\) with Haar law independently of a spectral
path with law \(Q^{\rm sp}\), and transport that path to physical
coordinates by \(U\).
For the filtration generated by \(U\) and observations through \(p\),
but not by \(X\), the exact likelihood from \cite[Equation~(2.9)]{FUA} is
\begin{equation}\label{edge:path-likelihood}
 \frac{\mathrm d\mathbf P}{\mathrm d\mathsf R}\bigg|_{\cF_p}
 =L_p(h_p),\qquad
 L_p(h)=\exp\{\phi_p(h)-\phi_p^{\rm sp}(h)\}.
\end{equation}
The notation \(L_p(h)\) suppresses its dependence on \(U\): in a Haar
average the spectral field \(U^{\mathsf T}h\) is held fixed, while
\(L_p(h)\) still varies with \(U\).  At every such field \(E_UL_p(h)=1\).
To check the likelihood,
the observation density is a common Gaussian factor times the posterior
partition function divided by the original one, and \(L_\circ\) cancels
the original partition functions.  In particular, for every nonnegative
spin or observation functional \(G\),
\begin{equation}\label{edge:bias-transfer}
 E_U[\ind_{\{L_\circ\ge1/2\}}Q_\mu G]\le2\mathbf E G.
\end{equation}

\begin{proposition}[Cap estimates with a tight lower threshold]
\label{prop:edge-estimates}
Fix \(0<C_1<\infty\), \(0\le R<\infty\), positive tolerances
\(\epsilon_{\rm sp},\epsilon_0,\epsilon_v,\epsilon_g\), and \(\zeta>0\).
For every \(\varepsilon>0\) there are finite positive constants
\(K,C,c,a,v_0\), a width \(s_0>0\), a radius \(R_{\rm ann}<\infty\),
and Gaussian disorder events
\(\mathcal C_n\) with
\(\liminf_n\Prob_{\rm GOE}(\mathcal C_n)\ge1-\varepsilon\), such that
the following hold on \(\mathcal C_n\), simultaneously at every scale
\begin{equation}\label{edge:scale-range}
 K\le k_p=np^3,\qquad p\le C_1r_n.
\end{equation}
\begin{enumerate}
\item The active dimension is between \(ck_p\) and \(Ck_p\),
 \(d_1\ge-\epsilon_{\rm sp}p^2\), and, with \(b_0=4/(3\pi)\),
 \begin{equation}\label{edge:trace-estimates}
 \left|\frac{n-\Tr R_{p,0}}{np}-b_0\right|\le\epsilon_{\rm sp},
 \qquad
 \left|\frac p n\Tr R_{p,0}^2-2b_0\right|\le\epsilon_{\rm sp}.
 \end{equation}
 On \(\|h\|\le Rp^2M_p\), the saddle satisfies
 \(-1+c_R\le\alpha_p(h)/p^2\le C_R\).  On \(|u|\le s_0p^2\),
 \(\ell_p(u)\asymp p^2M_p\) and \(\ell'_p(u)\asymp M_p\).

\item On the field ball \(\|h\|\le Rp^2M_p\),
 \begin{equation}\label{edge:static-estimates}
 \begin{split}
 K_p(h)&\preceq C_Rp^{-2}I_{\cH_p},\\
 |\phi_p(h)-\phi_p^{\rm sp}(h)|&\le\epsilon_v k_p,\\
 \|P_pm_p(h)-R_{p,\alpha_p(h)}h\|&\le\epsilon_gM_p .
 \end{split}
 \end{equation}
 The constant \(C_R\) in the first line may be chosen independently of
 the requested annular tolerance \(\epsilon_0\).
 On the full annulus \(|\alpha_p(h)|\le s_0p^2\), the field is nonzero;
 for \(e=h/\|h\|\),
 \begin{equation}\label{edge:coarse-curvature}
 K_p(h)\preceq(1+\epsilon_0)p^{-2}I_{\cH_p},
 \qquad p^2\langle e,K_p(h)e\rangle\le0.9 .
 \end{equation}
 The annulus lies in \(\|h\|\le R_{\rm ann}p^2M_p\), and the event also
 includes the static comparisons on this ball, with constants allowed
 to depend on \(R_{\rm ann}\).

\item For every \(v\ge v_0\), the uncapped law \(\mu\), every posterior
 \(\mu_{p,h}\) in the stated field ball, and their spherical counterparts
 satisfy
 \begin{equation}\label{edge:projection-tails}
 \Pr\{\|P_pX\|>vM_p\}\le C e^{-c k_pv^6}.
 \end{equation}
 The event is empty for \(v>p^{-1/2}\).  The spherical pair bound of
 Lemma~\ref{edge:spherical-tails} also holds for each number of replicas
 fixed in advance, with constants allowed to depend on that number.  Moreover,
 \begin{equation}\label{edge:ordinary-annulus}
 Q_\mu\{|\alpha_p(h_p)|>\zeta p^2\}\le C e^{-c k_p}.
 \end{equation}

\item For every scale in \eqref{edge:scale-range} with \(p\le r_n\),
 simultaneously for every \(0\le f\le1\),
 \begin{equation}\label{edge:variance-estimates}
 \begin{split}
 V_p(f)&\le Cp^{-2}\cE(f)+C e^{-c k_p^a},\\
 p^2Q_\mu\|\nabla_{\cH_p}F_p(h_p)\|^2
       &\le C V_p(f)+C e^{-c k_p^a}.
 \end{split}
 \end{equation}
 The event and constants in these inequalities are independent of \(f\).
\end{enumerate}
All constants in the saddle, length, and comparison bounds are
deterministic on \(\mathcal C_n\) and may depend on the fixed inputs and
the requested failure probability, but not on \(n\), \(p\), or the field
in its stated region.
\end{proposition}

Every radius, tolerance, and finite multiplicative enlargement of a scale
interval is fixed before the constants are chosen.  Increasing \(K\) by a
fixed factor supplies any such enlargement at the lower endpoint.
A finite number of further requests for cap estimates may be imposed by
intersecting their events, with a smaller failure probability assigned to
each.  This preserves the constants and exponent of any variance estimate
already obtained.

\paragraph{Reduction to positive covariance error.}
Fix \(W\) and \(0\le f\le1\).  In the time variable \(t=p^2\), let
\[
 \mathcal F_t^{\rm obs}=\sigma(h_q:0\le q\le\sqrt t),
 \qquad M_t^f=F_{\sqrt t}(h_{\sqrt t}).
\]
This filtration excludes the spin and includes any information in
\(h_0\); the initial precision \(B_0\) need not vanish.  The bounded
posterior martingale \(M^f\) is continuous.  On positive time intervals,
its predictable bracket has density
\(\|\nabla_{\cH_{\sqrt t}}F_{\sqrt t}(h_{\sqrt t})\|^2\).
The coefficient can change at an activation, but the martingale has no
jump there.  Total variance and the bracket identity therefore give
\begin{equation}\label{edge:variance-identity}
 -\frac{\mathrm dV_p(f)}{\mathrm d(p^2)}
       =Q_\mu\|\nabla_{\cH_p}F_p(h_p)\|^2
\end{equation}
almost everywhere in \(p>0\), and in integrated form across activation
points.  This is the localization variance calculation in
\cite[proof of Theorem~1, Equation~(14) and the display after Equation~(15)]{EldanKoehlerZeitouni2022},
with deterministic rate \(P_p\).

Let \(\mathcal N\) be a jointly measurable set of triples \((p,W,h)\)
with \(h\in\cH_p\), chosen independently of \(f\), and let
\(\eta_p=\eta_p(W)\ge0\) be a measurable scalar fixed at each
\((W,p)\).  Suppose that
\(p^2K_p(h)\preceq(1+\eta_p)I_{\cH_p}\) for \(h\in\cH_p\) when
\((p,W,h)\notin\mathcal N\).  Put
\begin{equation}\label{edge:common-failure-cost}
 Z_p(X)=\frac{\|P_pX\|^2}{M_p^2},\qquad
 \rho_p(W)=k_pQ_\mu\!\left[
       \ind_{\mathcal N}(p,W,h_p)\mu_{p,h_p}(Z_p)\right].
\end{equation}
Conditional covariance and Cauchy--Schwarz give, at the observed field,
\begin{equation}\label{edge:gradient-bounds}
 \begin{split}
 \|\nabla_{\cH_p}F_p(h_p)\|^2
       &\le\|K_p(h_p)\|\Var_{\mu_{p,h_p}}(f),\\
 p^2\|\nabla_{\cH_p}F_p(h_p)\|^2
       &\le k_p\mu_{p,h_p}(Z_p).
 \end{split}
\end{equation}
For the second bound, write
\(\nabla F_p=\mu_{p,h_p}[P_pX(f-F_p)]\) and use \(|f-F_p|\le1\).
Its right side is a common cost for every such \(f\).  Splitting the
expectation in \eqref{edge:variance-identity} at \(\mathcal N\) yields
\begin{equation}\label{edge:variance-differential}
 -\frac{\mathrm dV_p}{\mathrm d\log p}
       \le2(1+\eta_p)V_p+2\rho_p
\end{equation}
almost everywhere.  If \(\eta_u\le Ck_u^{-a}\) for a fixed \(a>0\) on
\([p,q]\), its backward integrating factor is at most
\[
 \exp\left(2\int_p^q(1+Ck_u^{-a})\,\mathrm d\log u\right)
       \le C(q/p)^2,\qquad
 \int_p^qk_u^{-a}\,\mathrm d\log u
       =\frac{k_p^{-a}-k_q^{-a}}{3a}\le\frac{K^{-a}}{3a}
\]
whenever \(k_p\ge K\).  In particular,
\begin{equation}\label{edge:variance-amplification}
 V_p\le C(q/p)^2V_q+
       C\int_p^q(u/p)^2\rho_u\,\mathrm d\log u .
\end{equation}
The remaining work for Item~4 is to obtain this covariance bound with
\(\eta_p=Ck_p^{-a}\), control the common cost \(\rho_p\) on dyadic
intervals, and use the source variance estimate at \(r_n\).  The
integrable relative error is what preserves the exponent \(2\).

\subsection{Spectral counts, saddles, and spherical tails}

For a specified spectral error, being small above a tight threshold means
\begin{equation}\label{edge:threshold-convention}
 \lim_{K\to\infty}\limsup_{n\to\infty}
 \Prob_{\rm GOE}\left\{
 \sup_{\substack{np^3\ge K\\p\le C_1r_n}}
       \text{the specified error at \(p\)}>\eta\right\}=0
 \quad(\eta>0).
\end{equation}
Only finitely many errors at a time are needed.  This is a double-limit
assertion, not convergence with \(n\) at a fixed \(K\).

\begin{lemma}[Spectral estimates above a tight threshold]
\label{edge:spectral}
In the convention \eqref{edge:threshold-convention}, the measures
\(\nu_p=k_p^{-1}\sum_i\delta_{d_i/p^2}\) converge on fixed bounded
intervals to \(\pi^{-1}\sqrt v\,\ind_{\{v>0\}}\mathrm dv\).  Their
distribution functions converge uniformly on compact intervals.
The same convention gives \(d_1/p^2\ge-o(1)\),
\(\dim\cH_p/k_p\to2/(3\pi)\), and the two limits in
\eqref{edge:trace-estimates}.  On spectral events of arbitrarily large
limiting probability there is also a fixed \(C\) such that
\begin{equation}\label{edge:counting-tail}
 \begin{aligned}
 \#\{i:d_i\le u\}&\le Cnu^{3/2}
       &&(u\ge cp^2,\ u\le4),\\
 \int_{[A,\infty)}(1+v)^{-2}\nu_p(\mathrm dv)&\le CA^{-1/2}
       &&(A\ge1),
 \end{aligned}
\end{equation}
for every fixed \(c>0\), after increasing the threshold.
Every spectral restriction below also includes eigenvalue-only events
\(\mathcal S_n^{\rm glob}\) of probability \(1-o(1)\) on which, uniformly
over their spectra,
\begin{equation}\label{edge:global-spectral}
 \frac1n\sum_i\delta_{\lambda_i}\ \weakto\
 \rho_{\rm sc}(\dd x):=\frac1{2\pi}\sqrt{4-x^2}\,
             \ind_{\{|x|\le2\}}\dd x,\qquad
 \lambda_1\longrightarrow2,\qquad \|W\|\le3 .
\end{equation}
\end{lemma}

\begin{proof}
We use the eigenvalue-only event of \cite[Proposition~3.1]{CS}.
For every deterministic \(L_n\to\infty\) with \(L_n\le\log n\), that
proposition gives positive denominators and, uniformly for
\(L_nn^{-2/3}\le u\le M\) with fixed \(M\),
\begin{equation}\label{edge:source-resolvents}
 \begin{split}
 m_1(u):=\frac1n\sum_i(d_i+u)^{-1}
  &=\frac{2+u-\sqrt{u(4+u)}}2+o(\sqrt u),\\
 m_2(u):=\frac1n\sum_i(d_i+u)^{-2}
  &\le C_Mu^{-1/2},\\
 \lambda_j&=2+o(L_nn^{-2/3})\quad(1\le j\le3),
 \qquad \|W\|\le3 .
 \end{split}
\end{equation}
The errors have deterministic uniform bounds on an event of probability
at least \(1-Ce^{-cL_n^{1/4}}\).

For \eqref{edge:global-spectral}, take \(\mathcal S_n^{\rm glob}\) to be
the separate source event with \(L_n=\sqrt{\log n}\).  Its probability
tends to one, and the edge and
norm assertions follow directly from \eqref{edge:source-resolvents}.
For every fixed \(u>1\), the first resolvent converges uniformly there
to
\[
 \frac{2+u-\sqrt{u(4+u)}}2
 =\int\frac{\rho_{\rm sc}(\dd x)}{2+u-x}.
\]
The norm bound makes the empirical spectral measures a compact family.
Every weak subsequential limit therefore has the displayed transform
for every \(u>1\), and uniqueness of the Stieltjes transform identifies
it with \(\rho_{\rm sc}\).  A contradiction by subsequences gives this
weak convergence uniformly over the source event.  Intersecting with
this event does not change the tight-threshold convention.

Let \(K_n\to\infty\) be arbitrary and set
\(L_n=\min\{K_n^{1/3},\sqrt{\log n}\}\), changing finitely many initial
values if necessary.  For \(np^3\ge K_n\),
\(L_nn^{-2/3}=o(p^2)\).  Thus, for every fixed \(v>0\),
\[
 m_1(vp^2)=1-\sqrt v\,p+O_v(p^2)+o_v(p)
\]
uniformly in the scale range, and \(d_1/p^2\ge-o(1)\).  For
\(u\ge cp^2\), the counted denominators satisfy \(0<d_i+u\le2u\).
Consequently
\[
 \#\{d_i\le u\}\le4u^2\sum_i(d_i+u)^{-2}\le Cnu^{3/2}.
\]
Summing on dyadic bands gives the second bound in
\eqref{edge:counting-tail}; the last band above \(u=4\) costs \(O(p)\)
by \(\|W\|\le3\), and obeys the same bound.

For fixed \(v,w>0\), the transform difference is exactly
\begin{equation}\label{edge:transform-difference}
 \int\left(\frac1{x+v}-\frac1{x+w}\right)\nu_p(\mathrm dx)
 =\frac{m_1(vp^2)-m_1(wp^2)}p
 \longrightarrow\sqrt w-\sqrt v.
\end{equation}
To identify the limit including the tail, put
\(\omega_p(\mathrm dx)=(1+x)^{-2}\nu_p(\mathrm dx)\).
These finite measures have bounded mass, the tail bound
\eqref{edge:counting-tail}, and negative support in an interval shrinking
to zero.  Every weak limit \(\omega\) is therefore supported on
\([0,\infty)\).  Taking \(w=1\) in \eqref{edge:transform-difference} gives
\[
 \int\frac{1+x}{x+v}\,\omega(\mathrm dx)=\frac1{1+\sqrt v}.
\]
Letting \(v\to1\) gives \(\omega([0,\infty))=1/2\), and hence
\[
 \int\frac{\omega(\mathrm dx)}{x+v}
       =\frac1{2(1+\sqrt v)^2}.
\]
Uniqueness of the Stieltjes transform identifies
\(\omega(\mathrm dx)=\pi^{-1}\sqrt x(1+x)^{-2}\mathrm dx\).  The
continuous distribution function of the corresponding limit of
\(\nu_p\) gives uniformity in endpoints on compact intervals.  Failure
of uniformity over \(p\) would provide a sequence of scales to which
this same compactness argument applies, a contradiction.

The clipped second trace has limit
\[
 \frac p n\Tr R_{p,0}^2
 =\int(\max\{x,1\})^{-2}\nu_p(\mathrm dx)
 \longrightarrow
 \frac1\pi\left(\int_0^1\sqrt x\,\mathrm dx+
                 \int_1^\infty x^{-3/2}\,\mathrm dx\right)
 =\frac8{3\pi}.
\]
For the first trace, subtract an ordinary resolvent before integrating:
\[
 \frac{n-\Tr R_{p,0}}{np}
 =\frac{1-m_1(p^2)}p+
 \int\left(\frac1{x+1}-\frac1{\max\{x,1\}}\right)\nu_p(\mathrm dx)
 \longrightarrow\frac4{3\pi}.
\]
The difference kernel is \(O(x^{-2})\).  Its limiting integral is
\(4/(3\pi)-1\), obtained by substituting \(x=t^2\) on each of
\((0,1)\) and \((1,\infty)\).

These assertions hold for every arbitrarily slowly divergent \(K_n\).
If \eqref{edge:threshold-convention} failed, one could choose
\(K_j\to\infty\) and \(n_j\to\infty\) with failure probability bounded
away from zero and embed \(K_{n_j}=K_j\) in a divergent deterministic
sequence.  The preceding uniform conclusion contradicts that subsequence.
\end{proof}

\begin{lemma}[Saddles and squared-norm densities]
\label{edge:norm-density}
On the spectral restrictions of Lemma~\ref{edge:spectral}, the saddle
and length conclusions in Proposition~\ref{prop:edge-estimates} hold.
For a saddle field in a fixed ball let
\[
 Y\sim N(m,\widehat R),\qquad \widehat R=R_{p,\alpha_p(h)},\quad m=\widehat Rh,\quad
 \sigma^2=2\Tr\widehat R^2+4m^{\mathsf T}\widehat Rm.
\]
Then \(\sigma^2\asymp n/p\), and the density \(f_Y\) of \(\|Y\|^2\)
satisfies \(f_Y(n)\asymp\sqrt{p/n}=p^2/\sqrt{k_p}\).
In standard-deviation units this density converges uniformly on
\(\mathbb R\) to the standard normal density as \(K\to\infty\).

For a fixed number \(m_0\) of replicas, suppose the spectral rows
\(Y_i\in\mathbb R^{m_0}\) are independent Gaussian vectors and at least
\(L>2\) row covariances obey
\(cu^{-1}I_{m_0}\preceq C_i\preceq Cu^{-1}I_{m_0}\).
Their means are arbitrary.  The joint density of
\(\mathcal Q=(\sum_iY_{i1}^2,\ldots,\sum_iY_{im_0}^2)\) obeys
\begin{equation}\label{edge:replica-norm-density}
 \sup_z f_{\mathcal Q}(z)\le C_{m_0}(u/\sqrt L)^{m_0}.
\end{equation}
This applies after rowwise Gaussian tilts which preserve the stated band
and independence of the rows.
\end{lemma}

\begin{proof}
Put \(t=u/p^2\) and
\(\mathfrak d_p(t)=(n-\Tr R_{p,tp^2})/(np)\).
For \(y=\|h\|/(p^2M_p)\), the saddle equation becomes
\begin{equation}\label{edge:scaled-saddle}
 \mathfrak d_p(t)=\frac{y^2}{(1+t)^2},\qquad t>-1,\qquad
 \mathfrak d'_p(t)=\frac p n\Tr R_{p,tp^2}^2.
\end{equation}
Lemma~\ref{edge:spectral} gives \(\mathfrak d_p(0)\) arbitrarily close
to \(b_0\), and
\[
 \mathfrak d'_p(t)\ge
       \frac{\dim\cH_p/k_p}{(1+t)^2}.
\]
It follows by integration that
\(\mathfrak d_p(-1+\kappa)<0\) for a sufficiently small fixed
\(\kappa>0\); no saddle can lie there or below.  On \(t\ge0\),
\(\mathfrak d_p(t)\ge b_0/2\), so \eqref{edge:scaled-saddle} bounds
\(t\) above when \(y\le R\).  The right side of \eqref{edge:saddle},
viewed as a function of its parameter, is strictly decreasing and
continuous, tends to infinity at \(-p^2\), and tends to zero at infinity.
This also proves existence and uniqueness.  Denominator comparison gives
\(\Tr\widehat R^2\asymp n/p\), \(\|m\|^2\le C_Rnp\), and
\(\|\widehat R\|\le C_Rp^{-2}\), hence \(\sigma^2\asymp n/p\).
For a sufficiently small fixed \(s_0\), both
\(\mathfrak d_p(t)\) and \(\mathfrak d'_p(t)\) are bounded above and
below on \(|t|\le s_0\).  Differentiating
\(\ell_p(tp^2)=p^2M_p(1+t)\sqrt{\mathfrak d_p(t)}\) proves the length
and derivative bounds.

We record the density argument used here and later.  For independent
\(Y_i=a_i+\sqrt{\upsilon_i}G_i\), suppose
\[
 \max_i\upsilon_i\le Cp^{-2},\qquad
 \sigma^2:=2\sum_i\upsilon_i^2+4\sum_i\upsilon_i a_i^2
       \asymp k_pp^{-4},
\]
and at least \(ck_p\) variances are at least \(cp^{-2}\).  The Gaussian
quadratic formula, as in \cite[Section~15.3, Equation~(15.5)]{CM}, gives
\begin{equation}\label{edge:norm-characteristic}
 \left|\E e^{it\sum_iY_i^2}\right|
 =\prod_i(1+4t^2\upsilon_i^2)^{-1/4}
   \exp\left(-\sum_i
       \frac{2t^2\upsilon_i a_i^2}{1+4t^2\upsilon_i^2}\right).
\end{equation}
Expansion of the centered log characteristic function through order two
gives
\begin{equation}\label{edge:norm-expansion}
 \log\E e^{it(\|Y\|^2-\E\|Y\|^2)/\sigma}
       =-t^2/2+O(k_p^{-1/2}|t|^3)
       \quad (|t|\le c\sqrt{k_p}).
\end{equation}
Indeed its remainder is at most
\[
 C|t/\sigma|^3\sum_i(\upsilon_i^3+\upsilon_i^2a_i^2)
       \le Ck_p^{-1/2}|t|^3.
\]
The band in \eqref{edge:norm-characteristic} bounds the modulus in these
units by \((1+ct^2/k_p)^{-c'k_p}\).  Its real logarithm is at most
\(-c_0t^2\) up to a small fixed multiple of \(\sqrt{k_p}\), and its
integral outside that interval is exponentially small in \(k_p\).
Fourier inversion and \eqref{edge:norm-expansion} give
\begin{equation}\label{edge:norm-local-limit}
 \sup_z\left|\sigma f_Y(\E\|Y\|^2+\sigma z)
       -(2\pi)^{-1/2}e^{-z^2/2}\right|\le Ck_p^{-1/2}.
\end{equation}
At the saddle \(\E\|Y\|^2=n\), proving the lower and upper density bounds.
At any mean, the same characteristic majorant gives the upper bound
\(Cp^2/\sqrt{k_p}\).

For replicas let \(T_t=\diag(t_1,\ldots,t_{m_0})\) and
\(H_i=C_i^{1/2}T_tC_i^{1/2}\).  The modulus of the row characteristic
function is at most \(\det(I+4H_i^2)^{-1/4}\); its extra factor from a
nonzero mean is at most one.  The singular values of \(H_i\) are at
least \(cu^{-1}\) times those of \(T_t\), including when the entries of
\(T_t\) have different signs.  Independence of the \(L\) rows bounds
the joint characteristic function by
\(\prod_j(1+ct_j^2/u^2)^{-L/4}\).  Its integral proves
\eqref{edge:replica-norm-density}.  This is the argument of
\cite[Section~15.3, Equation~(15.6)]{CM};
it explains the independence requirement on the spectral rows.
\end{proof}

\begin{lemma}[Spherical projection and pair tails]
\label{edge:spherical-tails}
Fix a field radius and a fixed number of replicas.  On the spectral
restrictions above, for every \(p\) with \(k_p\ge K\) and \(p\le C_1r_n\),
both the uncapped zero-field spherical law and the capped spherical law
at a field in that ball satisfy
\begin{equation}\label{edge:spherical-tail-bounds}
 \begin{aligned}
 \Pr\{\|P_pX\|>vM_p\}&\le C e^{-c k_pv^6}
       &&(v\ge v_0),\\
 \Pr\{|n^{-1}X^{\mathsf T}X'|>w\}&\le C e^{-c n w^3}
       &&(Cp\le w\le\rho_0).
 \end{aligned}
\end{equation}
Here \(X,X'\) are independent replicas at the same field,
\(\rho_0>0\) is fixed and small, and the second bound holds for every
pair of a fixed finite family.  The constants are uniform in \(p,n\).
\end{lemma}

\begin{proof}
We give the scale check in the polar argument of
\cite[Section~15.3, Equations~(15.5)--(15.8)]{CM}.
For an interaction \(A\), write \(R_{A,u}=(2+u-A)^{-1}\) and
\[
 q_{A,u}(h)=c_n+\tfrac n2(2+u)+\tfrac12\log\det R_{A,u}
                      +\tfrac12h^{\mathsf T}R_{A,u}h,\qquad
 c_n=\log\!\left(2^{n/2}\Gamma(n/2)n^{1-n/2}\right).
\]
Coarea in squared radius gives, exactly,
\begin{equation}\label{edge:frozen-density}
 \phi_A^{\rm sp}(h)=q_{A,u}(h)+\log f_{A,u,h}(n),\qquad
 \partial_uq_{A,u}(h)
   =\tfrac12\bigl(n-\Tr R_{A,u}-\|R_{A,u}h\|^2\bigr),
\end{equation}
where \(f_{A,u,h}\) is the density of
\(\|N(R_{A,u}h,R_{A,u})\|^2\).  In the capped case take \(A=W_p\)
and start at its saddle \(u=\alpha\).  For \(u\ge p^2\),
\[
 n-\Tr R_{p,u}
 \le Cn\sqrt u+Ck_pp^2/u^2\le Cn\sqrt u.
\]
The first term is \eqref{edge:source-resolvents}; the second compares
the clipped active denominators with \(d_i+u\), using
\(\dim\cH_p\le Ck_p\).  For \(u\) in a bounded range in \(p^2\) units,
the upper bound is \(C_Rnp\).  Integrating
\eqref{edge:frozen-density} yields
\begin{equation}\label{edge:frozen-normalizer-cost}
 q_{W_p,u}(h)-q_{W_p,\alpha}(h)\le Cnu^{3/2}+C_Rk_p,
\end{equation}
where the first constant is independent of the field radius.

For a rowwise quadratic test \(\mathcal T\) of \(m_0\) replicas and
\(\lambda>0\), let \(f_{\mathcal Q}^{u,\lambda}\) be the norm density
under the normalized Gaussian tilt \(e^{\lambda\mathcal T}\) of
\(m_0\) independent copies of \(N(R_{p,u}h,R_{p,u})\).  Exponential
Markov and polar coarea give
\begin{equation}\label{edge:polar-tail}
 \begin{split}
 (\mu_{p,h}^{\rm sp})^{\otimes m_0}\{\mathcal T\ge b\}
 \le{}&
 e^{m_0(q_{W_p,u}(h)-q_{W_p,\alpha}(h))-\lambda b}
 \E_u e^{\lambda\mathcal T}
 \frac{f_{\mathcal Q}^{u,\lambda}(n\mathbf1)}
      {f_{W_p,\alpha,h}(n)^{m_0}} .
 \end{split}
\end{equation}
This compares norm densities before and after the normalized tilt; it
does not divide by the probability of a shrinking shell.

For a pair take \(\mathcal T=\pm X^{\mathsf T}X'\), \(b=nw\),
\(u=\iota w^2\), and \(\lambda=c'u\), with \(c'>0\) a sufficiently
small absolute constant.  We will choose \(\iota>0\) small and then
the lower multiple of \(w/p\) large, so that \(u\gg p^2\).
The Gaussian formula gives
\begin{equation}\label{edge:quadratic-mgf}
 \log\E_u e^{\lambda\mathcal T}
 \le C\lambda^2\Tr R_{p,u}^2+C\lambda\|R_{p,u}h\|^2
 \le Cnu^{3/2}+C_Rk_p.
\end{equation}
Indeed \(\Tr R_{p,u}^2\le Cn/\sqrt u\) and
\(\|R_{p,u}h\|^2\le C_Rnp^5/u^2\); the pair quadratic has zero first
trace in the centered Gaussian formula.  For
\(\mathcal T=\|P_pX\|^2\), \(b=nw\), its extra first trace costs only
\(Ck_p\).  The choice of \(c'\) preserves positive precision and
covariance comparable to \(u^{-1}I_{m_0}\) on \(ck_p\) active rows.
Lemma~\ref{edge:norm-density} bounds the density ratio in
\eqref{edge:polar-tail} by \(C_{m_0}(u/p^2)^{m_0}\).

For uncapped zero field start at \(u=p^2\).  The resolvent gives
\(n-\Tr(d+p^2)^{-1}\asymp np\), a leading variance comparable to
\(p^{-2}\), and variance \(O(n/p)\) for the other squares.  Chebyshev's
inequality puts the residual squared length to be supplied by that
coordinate in \([cnp,Cnp]\) with probability \(1-C/k_p\).  On this
interval its Gaussian-square density is at least
\(c\sqrt{p/n}e^{-Ck_p}\).  Convolution gives the same lower density
scale, with this extra exponential factor.  The uncapped resolvent
also gives \eqref{edge:frozen-normalizer-cost}, and the tilted upper
density has the same band bound.

The logarithm of the resulting probability is at most
\[
 -c'\iota nw^3+C\iota^{3/2}nw^3+C_Rk_p
       +O_{m_0}(\log(u/p^2)).
\]
Choose \(\iota\) so that \(C\sqrt\iota<c'/4\), and then \(w/p\) larger
than a sufficiently large fixed constant.  Since
\(nw^3=k_p(w/p)^3\), the last two terms are absorbed by another
\(c'\iota nw^3/4\) when \(K\) is large.  This proves the pair bound.
For projection put \(w=v^2p\le1\); then \(nw^3=k_pv^6\).  For
\(v>p^{-1/2}\) the event is empty because \(\|X\|=\sqrt n\).
\end{proof}

\begin{corollary}[Growth of the limiting spectral rates]
\label{edge:limit-growth}
For GOE-almost every \(g\), and every \(b\) with \(b+g_1>0\), there are
finite positive constants \(c_{g,b},C_{g,b}\) such that
\begin{equation}\label{edge:limit-rate-bounds}
 c_{g,b}(1+i)^{2/3}\le g_i+b\le C_{g,b}(1+i)^{2/3}
 \qquad(i\ge1).
\end{equation}
They can be chosen measurably when \(b\) is a measurable admissible
function of \(g\).
\end{corollary}

\begin{proof}
Fix errors in Lemma~\ref{edge:spectral} for which the number of points
with \(d_i\le p^2\) is between \(c np^3\) and \(Cnp^3\), and apply it
at the countable scales \(np^3=2^jK\).  For any finite list, finite
edge-coordinate convergence transfers the inequalities, with slightly
enlarged constants, to counts of limiting points below
\((2^jK)^{2/3}\).  A fixed small buffer in the endpoints may be used;
the limiting process has no point at any of these countably many fixed
endpoints.  Increase the finite list and use continuity of probability
for decreasing events.  The uniform finite-\(n\) event shows that all
the limiting inequalities hold with probability arbitrarily close to one
when \(K\) is large.  Letting that failure probability tend to zero
shows that almost surely, after a finite threshold,
\[
 cL^{3/2}\le\#\{i:g_i\le L\}\le CL^{3/2}
\]
on dyadic \(L\), and hence on all large \(L\) by monotonicity.  Invert
these inequalities, add \(b\), and adjust the bounds for the finitely
many remaining positive numbers \(g_i+b\).  The least integer threshold
and rational bounds give measurable choices.
\end{proof}

\subsection{Comparison with the cube}

The spherical estimates control deterministic fields.  We next obtain
one disorder event on which the corresponding cube estimates hold for
every field and scale.  The additional work is confined to
\(k_p=O(\sqrt{\log n})\).

\begin{lemma}[Two-replica overlap cells]\label{edge:gram-cell}
Let \(\mathscr R_n=\{-1+2j/n:0\le j\le n\}\) and
\(I_r=[r-1/n,r+1/n]\).  The overlap of two independent uniform cube
spins has mass
\[
 p_n(r)=2^{-n}\binom{n}{n(1+r)/2},\qquad r\in\mathscr R_n,
\]
whereas the overlap of two independent uniform spherical spins has
density
\[
 f_n(s)=\frac{\Gamma(n/2)}{\sqrt\pi\,\Gamma((n-1)/2)}
             (1-s^2)^{(n-3)/2},\qquad -1<s<1.
\]
For a sufficiently small fixed \(\rho_0>0\), uniformly for large \(n\)
and \(r\in\mathscr R_n\) with \(|r|\le\rho_0\),
\begin{equation}\label{edge:pair-cell-mass}
 p_n(r)\le e^{C+Cnr^4}\int_{I_r}f_n(s)\,\mathrm ds.
\end{equation}
Let \(H_2\) be the uniform law of an ordered orthonormal pair
\((q,q')\).  If a positive measurable angular weight
\(a_n(s,q,q')\) has finite integrals and satisfies, for \(H_2\)-almost
every frame,
\[
 |\log a_n(s,q,q')-\log a_n(r,q,q')|\le K_0
 \quad\text{for almost every }s\in I_r,
\]
then
\begin{equation}\label{edge:pair-cell-weighted}
 p_n(r)\int a_n(r,q,q')\,\mathrm dH_2
 \le e^{C+Cnr^4+K_0}
       \int_{I_r}\!\int a_n(s,q,q')\,\mathrm dH_2\,f_n(s)\,\mathrm ds.
\end{equation}
The cells have disjoint interiors.  Summing over their midpoints in an
overlap interval enlarges the spherical interval by at most \(1/n\)
at each endpoint.
\end{lemma}

\begin{proof}
We derive this form for cells centered at lattice points from the two-replica calculations in
\cite[Section~15.4, Equation~(15.10)]{CM} and
\cite[proof of Lemma~3.2]{FUA}.  Set
\[
 I(r)=\tfrac12\bigl((1+r)\log(1+r)+(1-r)\log(1-r)\bigr),
 \qquad J(r)=-\tfrac12\log(1-r^2).
\]
Uniform Stirling bounds on \(|r|\le\rho_0\) put \(p_n(r)\) between
fixed multiples of \(n^{-1/2}e^{-nI(r)}\).  The same bounds for the
gamma ratio, and the bounded logarithmic oscillation of \(f_n\) on
\(I_r\), put its cell integral between fixed multiples of
\(n^{-1/2}e^{-nJ(r)}\).  Since \(J(r)-I(r)=O(r^4)\), this proves
\eqref{edge:pair-cell-mass}.  The assumed oscillation of the positive
angular weight gives \eqref{edge:pair-cell-weighted} by integration.
The final assertion follows directly from the cell endpoints.
\end{proof}

\begin{lemma}[A Haar second moment at the low scales]
\label{edge:replica-comparison}
On the spectral restrictions above, for every fixed spectral field in
a prescribed ball and every \(K\le k_p\), \(p\le C_1r_n\),
\begin{equation}\label{edge:likelihood-second-moment}
 E_U(L_p(h)-1)^2\le C_Rn^{-1/4}.
\end{equation}
The field is fixed inside the expectation; the constant is uniform over
the indicated fields and scales.
\end{lemma}

\begin{proof}
We adapt the two-replica calculation of
\cite[Lemma~3.2, Equations~(3.3)--(3.4)]{FUA}, whose stated observation
range starts at \(r_n\).  The estimate below the source range is proved
here.
Conditional on \(\rho=n^{-1}X^{\mathsf T}X'\), the Haar angular
integral for a cube pair is the spherical angular integral.  On the
ordered orthonormal sum-and-difference frame \((q,q')\), its exponent is
\begin{equation}\label{edge:pair-angular-exponent}
 \Psi_{p,h}(s;q,q')=
 \frac n2\bigl((1+s)q^{\mathsf T}W_pq
       +(1-s){q'}^{\mathsf T}W_pq'\bigr)
       +\sqrt{2n(1+s)}\,h^{\mathsf T}q.
\end{equation}
On a fixed \(|s|\le2\rho_0<1\), the spectral norm bound and the field
bound \(\|h\|\le R\sqrt n\,p^{5/2}\) give
\(|\partial_s\Psi_{p,h}|\le C_Rn\).  Thus the positive angular weight
\(e^{\Psi_{p,h}}\) has a fixed logarithmic oscillation on every
width-\(2/n\) cell, as required in \eqref{edge:pair-cell-weighted}.
Retain
\(|\rho|\le C'r_n\), choosing \(C'\) larger than the lower tail
constant in Lemma~\ref{edge:spherical-tails} times \(C_1\).
On dyadic bands between \(C'r_n\) and a small fixed \(\rho_0\),
the weighted cell comparison costs \(e^{C+Cn\rho^4}\), while
Lemma~\ref{edge:spherical-tails} gives \(Ce^{-cn\rho^3}\).
The spherical band endpoints move by at most \(1/n\), negligible
compared with \(r_n\).
Decreasing \(\rho_0\) makes the latter pay the former and leaves
\(Ce^{-cnr_n^3}\).  At overlaps bounded away from zero, the global
spectral law, edge limit, and norm bound in \eqref{edge:global-spectral}
supply the global inputs for the uncapped zero-field angular comparison
of \cite[Section~15.4]{CM}.  Indeed, for every fixed \(w>2\), they give
\[
 \frac1n\log\det(wI-W)\longrightarrow
       \int\log(w-x)\rho_{\rm sc}(\dd x);
\]
interlacing gives the same per-site limit for every orthogonal-hyperplane
compression.  The cited deficit-filling lower bound then gives the full
critical free-energy normalization by taking \(n\to\infty\) at fixed
\(w>2\), followed by \(w\downarrow2\).  For clarity, the normalized
zero-field contribution being bounded is
\[
 \mathcal M_n(\rho_0;\boldsymbol\lambda)
   =E_U\!\left[L_\circ^2\,
       \mu_W^{\otimes2}\{|n^{-1}X^{\mathsf T}X'|\ge\rho_0\}\right].
\]
Here \(\boldsymbol\lambda=(\lambda_1,\ldots,\lambda_n)\) is fixed,
\(W=U\diag(\boldsymbol\lambda)U^{\mathsf T}\), and \(\mu_W\) is the
zero-field cube Gibbs law for this \(W\).
The comparison in
\cite[Section~15.4, ``Zero field and overlaps bounded away from zero'']{CM},
together with the lower normalization in its Section~15.3, gives
\(\mathcal M_n(\rho_0;\boldsymbol\lambda)\le
\exp\{-\kappa_{\rho_0}n+o(n)\}\) for a fixed \(\kappa_{\rho_0}>0\),
along the retained spectral sequences with the preceding per-site
errors uniform.  This is a consequence of that proof; the norm bound
supplies continuity at the endpoints.
Changing to \(W_p,h\) changes weights and log normalizers by
\(O_R(np^2+np^{5/2})=o(n)\), so that region still contributes \(e^{-cn}\).

On retained cells of width \(2/n\), Stirling's formula for the binomial
overlap mass and the spherical density proportional to
\((1-\rho^2)^{(n-3)/2}\) gives relative error
\begin{equation}\label{edge:retained-stirling}
 O(nr_n^4+r_n+n^{-1}).
\end{equation}
This is the cell estimate in the cited replica proof, now with cutoff
\(C'r_n\) at every \(p\).  More precisely than the fixed bound above,
the oscillation of \eqref{edge:pair-angular-exponent} on one cell is at most
\[
 C\bigl(|q^{\mathsf T}W_pq-{q'}^{\mathsf T}W_pq'|+p^{5/2}\bigr),
\]
which is uniformly bounded.  The normalized spherical integral of
this expression on retained cells is \(O(r_n)\).  Indeed,
\[
 q^{\mathsf T}W_pq-{q'}^{\mathsf T}W_pq'
 =\frac{2X^{\mathsf T}W_pX'
       -\rho(X^{\mathsf T}W_pX+{X'}^{\mathsf T}W_pX')}
        {n(1-\rho^2)}.
\]
The diagonal terms cost \(O(r_n)\).  For the cross term, tilt a Gaussian
pair at the cap saddle by \(\pm cp^2X^{\mathsf T}W_pX'\), with \(c\)
small enough to preserve a fixed precision margin.  The Gaussian
logarithmic moment is \(O_R(k_p)\), since
\(\Tr R_{p,\alpha}^2=O(n/p)\) and
\(\|R_{p,\alpha}h\|^2=O_R(np)\).  The replica density bound
\eqref{edge:replica-norm-density} costs a fixed factor.  Hence
\[
 \Pr_{\rm sp}\{|X^{\mathsf T}W_pX'|>t\}
       \le C e^{C_Rk_p-cp^2t}.
\]
Integrating from a sufficiently large multiple of \(np\) gives
\(\E_{\rm sp}|X^{\mathsf T}W_pX'|\le C_Rnp\le C_Rnr_n\).
This full probability integral also bounds its unnormalized restriction
to retained cells.

The bounded oscillation permits the mean value theorem for exponentials.
After integration and summation over disjoint cells, the error is the
preceding integrated \(O(r_n)\); no supremum over frames is taken.
Slightly enlarged boundary cells are paid by the pair tail.  Thus
\[
 E_UL_p(h)^2
 \le1+C_R\{nr_n^4+r_n+n^{-1}+e^{-cnr_n^3}\}.
\]
Use \(E_UL_p(h)=1\),
\(nr_n^4=n^{-1/3+4\cdot10^{-4}}\), and
\(nr_n^3=n^{3\cdot10^{-4}}\) to obtain
\eqref{edge:likelihood-second-moment}.
\end{proof}

\begin{lemma}[Uniform cube cap estimates]\label{edge:cube-caps}
The static, curvature, and projection assertions
\eqref{edge:static-estimates}--\eqref{edge:projection-tails} hold on
ordinary Gaussian disorder events of arbitrarily large limiting
probability, with a tight threshold and every fixed field radius,
tolerance, and scale buffer prescribed in
Proposition~\ref{prop:edge-estimates}.  The projection bound also holds
for \(P_qX/M_q\) under a cap interpolated linearly between \(B_p\) and
\(B_q\), whenever \(q\ge p\), \(q/p\) is bounded, and the field lies
in a fixed \(q\)-field ball.
\end{lemma}

\begin{proof}
On dyadic \(p\) and \(v\), the spherical tail and
\eqref{edge:bias-transfer} give
\[
 E_U[\ind_{\{L_\circ\ge1/2\}}
       \mu\{\|P_pX\|>vM_p\}]
       \le C e^{-c k_pv^6}.
\]
Markov's inequality with half the exponent has total failure at most
\[
 \sum_{j,l\ge0}C e^{-cK\,8^j64^l}\le C e^{-c'K}.
\]
The discarded \(L_\circ<1/2\) event has probability \(o(1)\) by
\eqref{edge:haar-bias}.  Monotonicity of \(P_p\), and a fixed-factor
change in \(M_p\), extend the uncapped result to every \(p,v\).

The posterior tails follow deterministically.  For \(q\ge p\) of bounded
ratio, let \(B_\theta\) lie between \(B_p\) and \(B_q\).  It is positive,
supported on \(\cH_q\), and has norm at most \(Cq^2\).  Restricting the
uncapped normalizer to \(\|P_qX\|\le VM_q\) gives
\[
 \mu e^{-X^{\mathsf T}B_\theta X/2}\ge e^{-Ck_q}
\]
for a fixed large \(V\).  The zero-field capped law is even, so its
normalizer for a linear field is at least one.  On the shell
\(\|P_qX\|\asymp vM_q\), a field of norm at most \(Rq^2M_q\) costs
at most \(Rk_qv\).  This and the denominator cost are absorbed by
the uncapped \(e^{-c k_qv^6}\) tail for large \(v\).  This proves the
posterior tail, including the interpolated assertion.  The spherical
argument is identical.  The first two projection moments under all
these posteriors are consequently bounded by \(C_RM_q\) and \(C_RM_q^2\).

For nearby scales and fields, integrating first derivatives of the log
partition function gives
\begin{equation}\label{edge:potential-interpolation}
 |\Delta\phi|\le C_Rk_q\left(
       |\Delta\log p|+\frac{\|\Delta h\|}{q^2M_q}\right),
\end{equation}
and likewise for \(\phi^{\rm sp}\).  Indeed,
\(\|B_q-B_p\|\le Cq^2|\log(q/p)|\) on \(\cH_q\); its derivative is
controlled by the second projection moment and the field derivative
by the first.  Fields in the smaller space are compared in \(\cH_q\).

Cover \(K\le k_p\le C_0\sqrt{\log n}\), with fixed \(C_0\) large enough
to overlap the source cutoff.  Net scales from above and their spectral
field balls at resolution
\(\eta/(C_R(1+C_0\sqrt{\log n}))\) in the units of
\eqref{edge:potential-interpolation}.  Conditional on the spectrum,
these are deterministic nets.  Since \(\dim\cH_p\le Ck_p\), their
total size is
\[
 \exp\{O_{\eta,R}(\sqrt{\log n}\log\log n)\}=n^{o(1)}.
\]
At a net point \eqref{edge:likelihood-second-moment} and Chebyshev give
\(\Pr_U(|\log L_p(h)|>\eta/2)\le C_\eta n^{-1/4}\).
The union bound and \eqref{edge:potential-interpolation} yield
\begin{equation}\label{edge:absolute-value-comparison}
 \sup_{\substack{K\le k_p\le C_0\sqrt{\log n}\\
                  \|h\|\le Rp^2M_p}}
 |\phi_p(h)-\phi_p^{\rm sp}(h)|\le\eta
\end{equation}
with failure \(n^{-1/4+o(1)}\), in addition to the projection event.
For \(k_p\ge c\sqrt{\log n}\), the cap and interpolation results of
\cite[Proposition~2.6, Equation~(2.11), and Lemma~2.7]{FUA} apply as
stated, since \(C_1r_n\) is eventually below their fixed upper cutoff.
Together these prove the value comparison at every scale.

For the mean, use this comparison on a larger ball.  For unit
\(v\in\cH_p\), freeze \(\alpha=\alpha_p(h)\).  The polar formula and
Lemma~\ref{edge:norm-density} give
\[
 \log\mu_{p,h}^{\rm sp}
       e^{\pm t v^{\mathsf T}(X-R_{p,\alpha}h)}
       \le\tfrac12t^2v^{\mathsf T}R_{p,\alpha}v+C.
\]
The cube bound has the additional \(2\epsilon_vk_p\).
Jensen with \(t=\theta p^2M_p\), in the direction of the projected
mean discrepancy, gives
\[
 \frac{\|P_pm_p(h)-R_{p,\alpha}h\|}{M_p}
       \le C_R\theta+\frac{2\epsilon_v+C/k_p}{\theta}.
\]
Choose \(\theta\) from the required mean tolerance, then impose a smaller
value tolerance and increase \(K\).

It remains to prove the Hessian assertion on the low range.  Put
\(\widehat R=R_{p,\alpha_p(h)}\), \(m=\widehat Rh\), and
\(\sigma^2=2\Tr\widehat R^2+4m^{\mathsf T}\widehat Rm\).  For a unit \(v\in\cH_p\)
and \(|\theta|\le1\), the shift \(h+\theta pv\) fits a fixed larger
ball because \(p/(p^2M_p)=k_p^{-1/2}\).  Freezing \(\alpha\), the
exact polar formula is
\[
 \log\mu_{p,h}^{\rm sp} e^{\theta p v^{\mathsf T}(X-m)}
 =\tfrac12\theta^2p^2v^{\mathsf T}\widehat Rv+
       \log\frac{f_{\alpha,h+\theta pv}(n)}{f_{\alpha,h}(n)}.
\]
The shifted mean \(m_\theta=m+\theta p\widehat Rv\) satisfies
\[
 \frac{\sigma_\theta^2}{\sigma^2}=1+O_R(k_p^{-1/2}),\qquad
 \frac{n-\E_\theta\|Y\|^2}{\sigma_\theta}
       =-\frac{2\theta p\langle m,\widehat Rv\rangle}{\sigma}
           +O_R(k_p^{-1/2}).
\]
The leading term is uniformly bounded.  The density estimate
\eqref{edge:norm-local-limit}, at this argument and at zero, gives
\begin{equation}\label{edge:frozen-mgf-curvature}
 \begin{split}
 \log\mu_{p,h}^{\rm sp} e^{\theta p v^{\mathsf T}(X-m)}
   &=\tfrac12\theta^2v^{\mathsf T}\Gamma_{p,h}v+O_R(k_p^{-1/2}),\\
 \Gamma_{p,h}
   &=p^2\left(\widehat R|_{\cH_p}
       -\frac4{\sigma^2}(\widehat Rm)\otimes(\widehat Rm)\right).
 \end{split}
\end{equation}
This is uniform in \(h,v,\theta\).
Cauchy--Schwarz and \(\sigma^2\ge4m^{\mathsf T}\widehat Rm\) imply
\(0\preceq\Gamma_{p,h}\preceq C_RI\).
Equation~\eqref{edge:absolute-value-comparison} changes the log moment
by at most \(2\eta\), so the cube error is at most
\(\delta=C_RK^{-1/2}+2\eta\).

To justify differentiating this approximation, suppose a real \(Z\) obeys
\[
 \sup_{|\theta|\le1}
 \left|\log\E e^{\theta Z}-a\theta^2/2\right|\le\delta,\qquad
 0\le a\le C,\quad 0<\delta\le1/16.
\]
The values at \(\pm1\) bound \(\E e^{|Z|}\), and hence bound the
fourth derivative of \(\psi(\theta)=\log\E e^{\theta Z}\) on
\([-1/2,1/2]\).  Taylor's formula gives
\[
 \left|\psi''(0)
       -\frac{\psi(t)+\psi(-t)-2\psi(0)}{t^2}\right|\le C_Ct^2.
\]
The assumed error changes this quotient by at most \(2\delta/t^2\).
With \(t=\delta^{1/4}\), this proves
\(|\Var Z-a|\le C_C\sqrt\delta\).  Applying it to
\(Z=pv^{\mathsf T}(X-m)\), for every unit \(v\), yields
\begin{equation}\label{edge:hessian-approximation}
 \|p^2K_p(h)-\Gamma_{p,h}\|
       \le C_R\sqrt{C_RK^{-1/2}+2\eta}.
\end{equation}

On the annulus put \(t=\alpha/p^2\),
\(A=(n-\Tr\widehat R)/(np)\), \(B=(p/n)\Tr\widehat R^2\), and \(e=h/\|h\|\).
Since \(\widehat R|_{\cH_p}=p^{-2}(1+t)^{-1}I\) and \(m\) is parallel to \(e\),
\[
 \langle e,\Gamma_{p,h}e\rangle=\frac{B}{(1+t)B+2A}.
\]
The trace estimates and denominator comparison give
\(A=b_0+O(\epsilon_{\rm sp}+s_0)\) and
\(B=2b_0+O(\epsilon_{\rm sp}+s_0)\), so this is
\(1/2+O(\epsilon_{\rm sp}+s_0)\).  All other quadratic directions are
bounded above by \((1-s_0)^{-1}\).  Choose the spectral accuracy and
\(s_0\) small, then \(\eta\) small and \(K\) large in
\eqref{edge:hessian-approximation}.  This proves the annular curvature
and radial deficit.  On the entire ball, the saddle gap and the same
estimate give \(C_Rp^{-2}\), with \(C_R\) fixed before the annular
tolerance.  The source curvature assertion covers the higher range.
\end{proof}

\subsection{Location and covariance of the observation path}

The uniform cap event controls external fields.  We now locate the random
field produced by the Gibbs observation.  A shrinking annulus will give
an averaged covariance estimate; a fixed annulus gives the stronger
exponential tail needed in Proposition~\ref{prop:edge-estimates}.

We first state the additional eigenvalue-only restriction used for norm
conditioning.  Put
\[
 b_n=1+(-g_{1,n})_+,\qquad
 \ell_{a,n}=(g_{a,n}+b_n)^{-1},\qquad
 D_n=n^{1/3}-\sum_a\ell_{a,n},
\]
and let \(q_n\) be the density of
\(\sum_a\ell_{a,n}(\xi_a^2-1)\) for independent standard normal
\(\xi_a\).  For every \(\chi>0\), one may fix sufficiently large
integers \(m,J_\chi\) and positive constants \(c_\chi,C_\chi\) so that
an eigenvalue-only event \(\mathcal S_n^{\rm ann}(\chi)\) satisfies
\(\liminf_n\Prob(\mathcal S_n^{\rm ann}(\chi))\ge1-\chi\) and, on it,
\begin{equation}\label{edge:conditioned-spectral-event}
 \begin{gathered}
 b_n\le C_\chi,\qquad \ell_{a,n}\ge c_\chi\quad(1\le a\le m),\\
 g_{a,n}\ge C_\chi^{-1}a^{2/3}\quad(J_\chi<a\le n),\qquad
 q_n(D_n)\ge c_\chi.
 \end{gathered}
\end{equation}
These restrictions follow from the eigenvalue bounds in
\cite[Lemma~2.1]{FUA}, finite edge-coordinate convergence, and the
density argument in its Proposition~2.3.  We intersect this event with the local spectral
event of Lemma~\ref{edge:spectral} and \(\mathcal S_n^{\rm glob}\),
assigning a portion of the prescribed failure probability to it.  The
eigenvalue events of \cite[Lemma~2.5]{FUA} needed for
\eqref{edge:haar-bias} are retained separately.  All constants below
may depend on the prescribed confidence; no fixed density lower bound
on an event of probability \(1-o(1)\) is asserted.

\begin{lemma}[The observation annuli]\label{edge:path-annuli}
On the fixed spectral restrictions, there are \(a_{\rm an},c,C>0\)
such that for every dyadic interval \([q,2q]\) with fixed buffers
inside the enlarged scale range,
\begin{equation}\label{edge:shrinking-annulus}
 \mathbf P\left\{\exists p\in[q,2q]:
       |\alpha_p(h_p)|>Ck_p^{-1/10}p^2\right\}
       \le C e^{-c k_q^{a_{\rm an}}}.
\end{equation}
For every fixed \(\zeta>0\), after increasing the threshold,
\begin{equation}\label{edge:fixed-annulus-biased}
 \mathbf P\left\{\exists p\in[q,2q]:
       |\alpha_p(h_p)|>\zeta p^2\right\}
       \le C_\zeta e^{-c_\zeta k_q}.
\end{equation}
Consequently \eqref{edge:ordinary-annulus} holds simultaneously at
every scale on ordinary disorder events of arbitrarily large limiting
probability.
\end{lemma}

\begin{proof}
By \eqref{edge:haar-bias}, under \(\mathbf P\) the spectral spin is
spherical.  Its Gaussian representation is that of
\cite[Proposition~2.3]{FUA}.  The annulus calculation in
\cite[Lemma~3.1, Equation~(3.1)]{FUA} is stated for the source observation
range beginning at \(r_n\); the tight-threshold dyadic estimates below
extend that calculation to the scales used here.
Set \(D_i'=d_i+b_nn^{-2/3}=n^{-2/3}(g_{i,n}+b_n)\).
Before conditioning, let \(Y_i\) be independent \(N(0,(D_i')^{-1})\);
conditioning on \(\sum_iY_i^2=n\) gives the spherical law.
The definition of \(b_n\) gives \(\min_iD_i'\ge n^{-2/3}\), while
\eqref{edge:conditioned-spectral-event} supplies the fixed leading
variance lower bounds.  The density of
\(n^{-2/3}\sum_iY_i^2\) at \(n^{1/3}\) is exactly \(q_n(D_n)\),
and hence is bounded below on the chosen spectral set.  The same set
and the active count give
\begin{equation}\label{edge:leading-resolvent-sum}
 \sum_{d_i\le p^2}(D_i')^{-1}
   =n^{2/3}\sum_{d_i\le p^2}\ell_{i,n}
   \le Cn^{2/3}(1+k_p^{1/3})\le Cnp.
\end{equation}
Indeed the first \(J_\chi\) terms are bounded by one, and the remaining
terms are bounded by \(C_\chi i^{-2/3}\).

Let \(\vartheta=p^2+b_nn^{-2/3}\) and, before conditioning, write
\(h_{p,i}=B_{p,i}Y_i+B_{p,i}^{1/2}Z_i\).
The quadratic form obtained from
\(\|h_p\|^2-\vartheta^2\|Y\|^2\) has mean
\(-\vartheta^2\Tr R_{p,b_nn^{-2/3}}\).  On an active coordinate the
coefficient of the standardized \(Y_i^2\) term is
\[
 \frac{B_{p,i}^2-\vartheta^2}{D_i'}=D_i'-2\vartheta=O(p^2);
\]
the cross coefficient is \(O(p^3/\sqrt{D_i'})\).
Equation~\eqref{edge:leading-resolvent-sum} and the capped second trace
bound therefore show that the matrix of this quadratic form, in the
standardized \((Y,Z)\) variables, satisfies
\[
 \|\mathsf A_p\|_{\rm F}^2\le Cp^4k_p,\qquad
 \|\mathsf A_p\|\le Cp^2k_p^{1/3}.
\]
The Gaussian quadratic-form moment formula for its centered version bounds a deviation of
\(p^2k_p^{9/10}\) by
\[
 C\exp\left[-c\min\{k_p^{4/5},k_p^{17/30}\}\right].
\]
This remains true after norm conditioning, with a constant loss.
Indeed use the same exponential tilt, with parameter at most one
eighth of \(\|\mathsf A_p\|^{-1}\).  Its joint covariance and its
\(Y\)-marginal covariance stay comparable to the originals.  A fixed
number of the leading marginal variances bounds the density of
\(n^{-2/3}\|Y\|^2\) from above by a constant, through the characteristic
formula \eqref{edge:norm-characteristic}.  The original density at
\(n^{1/3}\) is bounded below on the chosen spectral set.  The polar
density ratio therefore costs a constant.

After conditioning, the deterministic target for \(\|h_p\|^2\) is
\(\ell_p(b_nn^{-2/3})^2\asymp p^2k_p\).  The preceding error is a
relative \(O(k_p^{-1/10})\).  It first places the inferred saddle in
a fixed annulus.  On that annulus
\(\mathrm d(\ell_p(u)^2)/\mathrm du\asymp k_p\).  Hence
\[
 |\alpha_p(h_p)-b_nn^{-2/3}|\le Cp^2k_p^{-1/10}.
\]
Since \(b_nn^{-2/3}/p^2=O(k_p^{-2/3})\), this proves the fixed-scale
version of \eqref{edge:shrinking-annulus}.

For a whole interval use a mesh of spacing \(k_q^{-3}\) in \(\log p\).
Its size is polynomial in \(k_q\).  Except with \(Ce^{-ck_q}\)
probability, the spherical projection at the upper endpoint is at most
\(CM_q\), by Lemma~\ref{edge:spherical-tails}.  The signal increments
on a mesh gap are then at most \(Cq^2M_qk_q^{-3}\).
The noise increments have covariance norm \(O(q^2k_q^{-3})\).
A constant-resolution net of the sphere in \(O(k_q)\) dimensions and
the Gaussian maximal inequality show that their norm on every gap is
less than \(q^2M_qk_q^{-1/10}\), except with \(Ce^{-ck_q}\)
probability.  The target length is Lipschitz in the same scaled units:
between activations differentiate the capped trace, using the active
count and second trace; the target is continuous at activations.
These observations interpolate the fixed-scale bound.  The polynomial
mesh size is absorbed by decreasing the positive exponent, proving
\eqref{edge:shrinking-annulus}.

For a fixed \(\zeta\), it suffices by inclusion to treat \(\zeta\)
below a fixed small annulus width.  The stronger exponent follows from a
different Gaussian reference, following the source-range fixed saddle
test in \cite[Section~19.2, Equation~(19.9) and its proof]{CM}.
Fix a small \(\xi>0\), to be chosen after \(\zeta\), and let
\[
 R_\xi=(2+\xi p^2-W)^{-1},\qquad
 Y\sim N(0,R_\xi),\qquad h_p=B_pY+\beta_p .
\]
Increase \(K\) so that \(d_1\ge-\xi p^2/2\).
The vector \(a_{p,\xi}=R_{p,\xi p^2}h_p\) is centered Gaussian with
covariance
\[
 C_{p,\xi}=R_\xi-R_{p,\xi p^2}.
\]
The resolvent and clipped traces give
\[
 \left|\Tr C_{p,\xi}-(n-\Tr R_{p,0})\right|
       \le C\sqrt\xi\,np+o_K(np),\qquad
 \|C_{p,\xi}\|\le C(\xi p^2)^{-1},\qquad
 \Tr C_{p,\xi}^2\le Cn/(\sqrt\xi\,p).
\]
For a fixed multiple of \(np\), exponential Markov with both signs
of a sufficiently small multiple of \(\xi p^2\) therefore gives
a deviation cost \(e^{-c_\zeta\xi k_p}\), once \(\xi\) is sufficiently
small.  The centered logarithmic moment costs \(O(\xi^{3/2}k_p)\),
whereas the Markov gain is a fixed multiple of \(\zeta\xi k_p\).

We spell out the norm-conditioning cost.  The deficit
\(n-\Tr R_\xi\) is between fixed multiples of \(\sqrt\xi\,np\),
after increasing \(K\) for this fixed \(\xi\).  A leading variance is
comparable to \((\xi p^2)^{-1}\); the other squared norms have variance
at most \(Cn/(\sqrt\xi\,p)\).  Chebyshev places the residual length
in a fixed interval comparable to \(\sqrt\xi\,np\) when
\(K\gg\xi^{-3/2}\).  On this interval the leading Gaussian-square
density is at least
\[
 C_\xi^{-1}\sqrt{p/n}\exp(-C\xi^{3/2}k_p).
\]
This is a lower bound for the original norm density by convolution.
Under the quadratic tilt used above, the marginal covariance of \(Y\)
retains \(ck_p\) independent spectral rows with covariance comparable
to \(p^{-2}\); the tilt keeps the joint covariance comparable.
Equation~\eqref{edge:replica-norm-density} bounds its norm density by
\(C_\xi\sqrt{p/n}\).  Thus conditioning costs at most
\(C_\xi e^{C\xi^{3/2}k_p}\).

Since \(R_{p,0}h_p=(1+\xi)R_{p,\xi p^2}h_p\), compare the saddle
length at the two boundaries \(\alpha=\pm\zeta p^2\).  The derivative
bounds in \eqref{edge:scaled-saddle}, and monotonicity beyond those
boundaries, put
\(\{|\alpha_p(h_p)|>\zeta p^2\}\) inside the relaxed squared-length
deviation just bounded.  Choose \(\xi\) so that its \(O(\sqrt\xi)\)
target error fits the \(\zeta\) tolerance and
\(C\xi^{3/2}<c_\zeta\xi/2\).  The conditioned deviation then has
probability \(C_\zeta e^{-c_\zeta'k_p}\).
The same \(k_q^{-3}\) mesh proves \eqref{edge:fixed-annulus-biased}.

Finally take \(q_j=2^j(K/n)^{1/3}\).  By
\eqref{edge:bias-transfer} and Markov, the ordinary whole-interval
probability is at most \(Ce^{-c k_{q_j}}\) outside disorder probability
\(Ce^{-c k_{q_j}}\), after decreasing \(c\).  Their sum is
\[
 \sum_{j\ge0}Ce^{-cK8^j}\le Ce^{-c'K}.
\]
Use \eqref{edge:haar-bias} for \(L_\circ<1/2\) and increase \(K\).
Every scale lies in one such interval and \(k_p\asymp k_{q_j}\)
there, proving \eqref{edge:ordinary-annulus} simultaneously in \(p\).
\end{proof}

We now fix the finite set of auxiliary radii, buffers, and tolerances
used in the covariance argument.  Fix once a sufficiently small
\(\gamma>0\), as specified by the compression argument below.  Then
choose finite \(C_*\) so that the range through \(C_*r_n\) contains
all parent intervals \([q/(2\gamma),q/\gamma]\), \(q\le r_n\), and
their prescribed enlargements, as well as the range through \(C_1r_n\).
Choose \(\eta\) sufficiently small compared with \(\gamma^2\), and an
auxiliary curvature tolerance \(\epsilon_{\rm aux}\) sufficiently small
compared with \(\eta\); it may be smaller than the prescribed
\(\epsilon_0\).  The extension construction below prescribes an
auxiliary width \(s_{\rm aux}\), static tolerances, radii, and a threshold
for these choices.  They are further cap requests under the
finite-intersection convention after Proposition~\ref{prop:edge-estimates}.
Let \(\mathcal G_n^{\rm aux}\) be the event, measurable from \(W\), on which
Lemma~\ref{edge:cube-caps} holds with these choices.  Let
\(\mathcal G_n^{\rm FUA}\) be the entire good-disorder event
\(\mathcal G_n\) defined in \cite[Section~3]{FUA} for its
Proposition~4.2, with that source's fixed tolerances and full scale and
field ranges.  This event has ordinary Gaussian-disorder probability
tending to one.  Set
\(\mathcal G=\mathcal G_n^{\rm aux}\cap\mathcal G_n^{\rm FUA}\).
Below we use the source covariance conclusion on the terminal interval
\([r_n,C_\gamma r_n]\), which is contained in its stated interval
\([r_n,p_*]\) for all sufficiently large \(n\).
All probabilities below are measured under \(\mathbf P\)
before restriction to \(\mathcal G\).  On each chosen spectral set,
the constants are uniform.

\begin{lemma}[Integrated score and covariance error]
\label{edge:score-covariance}
Let \(I=[q,2q]\), with a prescribed fixed multiplicative enlargement
inside the scale range.  At path fields write
\[
 a_p=P_p(m_p-m_p^{\rm sp}),\qquad D_p=p^2K_p-I_{\cH_p}.
\]
There are \(c,C,a_1>0\) such that
\begin{equation}\label{edge:score-bound}
 \mathbf P\left\{
   \int_I\|a_p\|^2\,\mathrm d(p^2)>Ck_q^{1/2}\right\}
       \le C e^{-c k_q^{a_1}},
\end{equation}
and
\begin{equation}\label{edge:frobenius-bound}
 \mathbf P\left(
 \mathcal G\cap
 \left\{\int_q^{2q}\frac{\|D_p\|_{\rm F}^2}{k_p}\,\mathrm d\log p
              >Ck_q^{-1/5}\right\}\right)
       \le C e^{-c k_q^{a_1}} .
\end{equation}
The same statements hold for every fixed enlargement of \(I\), with
adjusted constants.  They are probability statements; no expectation
of the untruncated score on their failures is asserted.
\end{lemma}

\begin{proof}
Put \(t=p^2\), and include \(U\) in the observation filtration, so that
every disorder event is measurable at time zero.  The exact likelihood
identity \eqref{edge:path-likelihood} yields
\begin{equation}\label{edge:likelihood-sde}
 \mathrm d\log L_{\sqrt t}
       =a_{\sqrt t}\cdot\mathrm d\mathcal B_t
          +\tfrac12\|a_{\sqrt t}\|^2\,\mathrm dt,
\end{equation}
where \(\mathcal B\) is the continuous cube innovation martingale with
\(\mathrm d\langle\mathcal B\rangle_t=P_{\sqrt t}\,\mathrm dt\);
the coefficients lie in the active space.  This is the exact filtering
calculation in \cite[proof of Lemma~3.5, after Equation~(3.7)]{FUA};
it follows as well by differentiating the posterior partition function,
including the interaction change \(-\mathrm dB_{\sqrt t}\).  There is no jump
in \(\log L\) at activation.

At a deterministic endpoint and for \(H>0\),
\eqref{edge:path-likelihood} gives
\[
 \mathbf P(\log L_p<-H)\le e^{-H}.
\]
For a fixed \(R_0\) containing the path annulus, put
\(\mathcal U_p=\{\widetilde h:\|\widetilde h\|\le R_0p^2M_p\}\).
Lemma~\ref{edge:replica-comparison} gives
\[
 \mathbf P(\log L_p>H,h_p\in\mathcal U_p)
 \le e^{-H}\int_{\mathcal U_p}Q^{\rm sp}(\mathrm d\widetilde h)\,
          E_U L_p(\widetilde h)^2
 \le (1+Cn^{-1/4})e^{-H}.
\]
The field in the inner expectation is deterministic.
Lemma~\ref{edge:path-annuli} excludes the complementary endpoint fields
with probability \(Ce^{-c k_q^{a_{\rm an}}}\).

On an enlarged interval write
\(\mathcal J_I=\int\|a_p\|^2\,\mathrm d(p^2)\) and
\(M=\int a_p\cdot\mathrm d\mathcal B_{p^2}\).
On the two endpoint events just given,
\(M+\mathcal J_I/2\le2H\).  If \(\mathcal J_I\ge8H\), then
\(M\le-\mathcal J_I/4\).
The positive exponential supermartingale
\(\exp(-M/4-\mathcal J_I/32)\) bounds this event by \(e^{-H/4}\).
Taking \(H=Ck_q^{1/2}\) proves \eqref{edge:score-bound}.

Between activations the two filtering equations give, with
\(\Delta_t=K_{\sqrt t}-K_{\sqrt t}^{\rm sp}\),
\[
 \mathrm da_{\sqrt t}
       =\Delta_t\,\mathrm d\mathcal B_t
          -K_{\sqrt t}^{\rm sp}a_{\sqrt t}\,\mathrm dt .
\]
Choose a smooth nonnegative cutoff \(\chi(t)\), equal to one on the
target interval, zero at the ends of a fixed enlargement, and with
\(|t\chi'(t)|\le C\).  The integrated Itô identity is
\begin{equation}\label{edge:score-energy}
 \begin{split}
 \int\chi t\|\Delta_t\|_{\rm F}^2\,\mathrm dt
 ={}&-\int(\chi+t\chi')\|a_{\sqrt t}\|^2\,\mathrm dt
       +2\int\chi t\langle a_{\sqrt t},
                         K_{\sqrt t}^{\rm sp}a_{\sqrt t}\rangle\,\mathrm dt\\
 &-\mathcal M
   -\sum_j\chi(t_j)t_j
       \|\Delta P_j(m_{\sqrt{t_j}}-m_{\sqrt{t_j}}^{\rm sp})\|^2,\\
 \mathcal M={}&2\int\chi t
       \langle a_{\sqrt t},\Delta_t\,\mathrm d\mathcal B_t\rangle .
 \end{split}
\end{equation}
Here \(\Delta P_j\) is the jump of \(P_{\sqrt t}\) at the activation
time \(t_j\).
The sum has the favorable sign.  On \(\mathcal G\) and the whole-path
annulus event, the coarse bounds give
\(\|tK_{\sqrt t}\|+\|tK_{\sqrt t}^{\rm sp}\|\le C\).
Thus the left side on the target interval is bounded by
\(C\mathcal J_I+C|\mathcal M|\).

For concentration, stop only the stochastic integral at the first
annulus failure, and set its integrand to zero on \(\mathcal G^c\).
The latter event is measurable at time zero.  The stopped bracket is
at most \(C\mathcal J_I\), and the stopped integral agrees with \(\mathcal M\)
on \(\mathcal G\) and the whole-path annulus event.  The exponential
bracket inequality, \eqref{edge:score-bound}, and
\eqref{edge:shrinking-annulus} consequently bound the target integral
of \(\|p^2(K_p-K_p^{\rm sp})\|_{\rm F}^2\,\mathrm d\log p\) by
\(Ck_q^{1/2}\), with the required exceptional probability intersected
with \(\mathcal G\).  The full identity \eqref{edge:score-energy} was
not stopped, so it has no terminal boundary term.

On the shrinking annulus, rotational invariance inside the flat cap
gives one spherical transverse covariance eigenvalue
\(\kappa_\perp\) and one radial eigenvalue.  The frozen linear moment
bound at shifts \(\pm pe\) gives
\[
 |e^{\mathsf T}(m_p^{\rm sp}-R_{p,\alpha}h_p)|\le C/p.
\]
Since \(\kappa_\perp=(e^{\mathsf T}m_p^{\rm sp})/\|h_p\|\) and
\(\|h_p\|\asymp p^2M_p\),
\[
 |p^2\kappa_\perp-1|
       \le C(k_p^{-1/10}+k_p^{-1/2}).
\]
The scaled radial covariance is bounded by the spherical version of
\eqref{edge:hessian-approximation}.  There are \(O(k_p)\) transverse
dimensions and one radial dimension, so
\[
 \frac{\|p^2K_p^{\rm sp}-I\|_{\rm F}^2}{k_p}
       \le C(k_p^{-1/5}+k_p^{-1}).
\]
Combine this with the preceding cube--sphere covariance estimate,
divide by \(k_q\asymp k_p\), and decrease \(a_1>0\) if necessary.
This proves \eqref{edge:frobenius-bound}.
\end{proof}

\subsection{The covariance iteration and exact-scale variance}

The Frobenius estimate says that most directions have a small covariance
error on average over a scale interval.  To control the largest positive
error, compare a child cap to a larger parent cap.  The radial deficit
puts every positive error in the tangent space of the parent field;
the unused Haar rotation then dilutes it in the smaller child space.

Use the fixed small \(\gamma>0\) above.  Let \(q\) be a child scale and \(p\) a parent
scale with \(x=q/p\in[\gamma,2\gamma]\), and put \(\Delta=B_p-B_q\).
Conditional on \(W\) and observations through \(q\), \(h_p\) has law
\begin{equation}\label{edge:transition}
 \mathsf T_{q,p,h_q}(\mathrm dy)
 =\frac{e^{\phi_p(y)}N(h_q,\Delta)(\mathrm dy)}
        {\int e^{\phi_p(z)}N(h_q,\Delta)(\mathrm dz)} .
\end{equation}
The Gaussian is on the positive-increment space.  Its covariance
compressed to \(\cH_q\) is
\(\Delta_q=(1-x^2)p^2I_{\cH_q}\).  Differentiation of the convolution
gives the exact identity from \cite[Section~4.1, Equation~(4.3)]{FUA},
\begin{equation}\label{edge:transition-covariance}
 K_q=\Delta_q^{-1}
       [\Cov_{\mathsf T_{q,p,h_q}}(h_p)]_{\cH_q}\Delta_q^{-1}
       -\Delta_q^{-1}.
\end{equation}

\begin{lemma}[An extension on the parent annulus]
\label{edge:annular-extension}
Choose \(\eta>0\) sufficiently small compared with \(\gamma^2\), then
the coarse curvature tolerance \(\epsilon_{\rm aux}\) sufficiently small
compared with \(\eta\).  One can choose \(s_{\rm aux}\), the static tolerances,
and the threshold \(K\) so that the following holds on \(\mathcal G\).
For every child field satisfying \(|\alpha_q(h_q)|\le s_{\rm aux}q^2\), put
\[
 \mathcal A_p=\{h:|\alpha_p(h)|\le s_{\rm aux}p^2/2\},\qquad
 \mathcal A_p^+=\{h:|\alpha_p(h)|<3s_{\rm aux}p^2/4\}.
\]
There is a semiconcave extension \(\bar\phi_p\) with upper Hessian bound
\((1+\eta)p^{-2}I\), agreeing with \(\phi_p\) on \(\mathcal A_p^+\).
Moreover,
\(\dist(\mathcal A_p,(\mathcal A_p^+)^c)\ge c_{\rm buf}p^2M_p\)
for a fixed \(c_{\rm buf}>0\).  Under both \(\mathsf T_{q,p,h_q}\)
and the transition obtained by replacing \(\phi_p\) with \(\bar\phi_p\),
the probability of \(\mathcal A_p^c\) is at most \(Ce^{-c_1k_p}\).
The same bound holds with any fixed polynomial
moment of \(\|h_p\|/(p^2M_p)\) inserted, after decreasing \(c_1>0\).
All constants are uniform in the indicated fields and scales.
\end{lemma}

\begin{proof}
We verify the deterministic extension proof of \cite[Lemma~4.3]{FUA},
following \cite[Proposition~16.1]{CM}, with the new lower threshold.
On \(\{|\alpha_p|\le3s_{\rm aux}p^2/4\}\) take the infimum, over \(z\), of
\[
 \phi_p(z)+\nabla\phi_p(z)\cdot(y-z)
       +\frac{1+\eta}{2p^2}\|y-z\|^2.
\]
Choose \(s_{\rm aux}\ll\eta,\gamma^2\).  The length bounds in
Lemma~\ref{edge:norm-density} give a fixed
\(\beta_{\rm buf}>0\) such that a segment of length at most
\(\beta_{\rm buf}p^2M_p\) starting in this smaller annulus remains in
the full annulus.  For such pairs the supporting inequality follows
by integrating the coarse Hessian bound.  For
\(\|y-z\|\ge\beta_{\rm buf}p^2M_p\), freeze the spherical saddle at
\(z\).  The polar linear-moment bound gives the spherical quadratic
support with Hessian \((p^2+\alpha_p(z))^{-1}I\), its saddle gradient,
and a bounded norm-density error.  The static value and mean estimates
add \(2\epsilon_vk_p+\epsilon_gM_p\|y-z\|\).  Relative to
\(\|y-z\|^2/(2p^2)\), the total residual is bounded by
\[
 Cs_{\rm aux}+\frac{4\epsilon_v+C/k_p}{\beta_{\rm buf}^2}
       +\frac{2\epsilon_g}{\beta_{\rm buf}} .
\]
Choose \(s_{\rm aux}\), then \(\beta_{\rm buf}\), then \(\epsilon_v,\epsilon_g\)
so that this is less than \(\eta\), and finally increase \(K\).
Every support lies above \(\phi_p\) on the smaller annulus.  Its
infimum is finite and semiconcave with the stated upper Hessian, and
it agrees there by choosing \(z=y\).  The length derivative gives
\(\dist(\mathcal A_p,(\mathcal A_p^+)^c)\ge c_{\rm buf}p^2M_p\).

For the transition tails freeze the child saddle
\(\alpha=\alpha_q(h_q)\), and use as reference the Gaussian transition
tilted by the parent frozen quadratic \(q_{W_p,\alpha}\).
Writing \(R_u=R_{u,\alpha}\), Gaussian convolution gives
\begin{equation}\label{edge:reference-transition}
 R_ph_p\sim N(R_qh_q,R_q-R_p),\qquad R_q^{-1}=R_p^{-1}-\Delta.
\end{equation}
The child saddle equation shows that
\(\E\|h_p\|^2=\ell_p(\alpha)^2\), exactly.  The field covariance has
norm at most \(C\gamma^{-2}p^2\).  Gaussian concentration of its norm
therefore gives
\[
 \Pr_{\rm ref}\{|\|h_p\|-\ell_p(\alpha)|>vp^2M_p\}
       \le C e^{-c\gamma^2k_pv^2},
 \qquad v\ge C/(\gamma\sqrt{k_p}).
\]
The reference radius lies strictly inside \(\mathcal A_p\):
\(|\alpha|\le s_{\rm aux}q^2\le4s_{\rm aux}\gamma^2p^2\).
The length derivative gives a fixed positive gap in units \(p^2M_p\)
to the boundary of \(\mathcal A_p\).

Choose \(z_0\in\cH_p\) with \(\|z_0\|=\ell_p(\alpha)\) and calibrate the
reference quadratic by
\[
 q^\sharp(y)=q_{W_p,\alpha}(y)+\log f_{W_p,\alpha,z_0}(n).
\]
The active covariance is scalar, so the added density depends only on
\(\|z_0\|\); the additive constant leaves the reference transition
unchanged.  On bounded scaled radii, the noncentral upper norm-density
bound and the static value comparison give
\(\phi_p(y)-q^\sharp(y)\le\epsilon k_p+O(1)\), with any prescribed fixed
\(\epsilon>0\).
Outside a sufficiently large bounded radius, the zero-field projection
tail gives
\[
 \phi_p(y)-\phi_p(0)
 \le Ck_p\left(1+
       \left(\frac{\|y\|}{p^2M_p}\right)^{6/5}\right).
\]
This follows by maximizing the exponent \(uv-cv^6\) after integrating
the tail against a linear field of scaled length \(u\).  Its growth is
strictly subquadratic.  For the extension choose the support centered
at the point of radius \(\ell_p(\alpha)\) in the direction of \(y\).
At distance \(vp^2M_p\), its excess over the reference is at most
\[
 \epsilon k_p(1+v)+C(\eta+s_{\rm aux})k_pv^2+O(1).
\]
This uses the static mean error for the linear term and the difference
of the two quadratic coefficients for the last term.  Choose
\(\eta+s_{\rm aux}\ll\gamma^2\), and then \(\epsilon\) small compared with
the squared fixed gap to \(\mathcal A_p^c\).  Outside \(\mathcal A_p\) these errors
are smaller than the reference Gaussian deviation cost.

The normalizing integrals relative to the reference are at least
\(e^{-\epsilon'k_p}\) for any needed fixed \(\epsilon'>0\).
To see this, integrate over a thin fixed shell about the reference radius,
contained in \(\mathcal A_p\).  Its reference probability is bounded below
when \(K\) is large.  On that shell the frozen density formula
\eqref{edge:frozen-density} bounds the off-saddle logarithmic density
loss by \(C\tau^2k_p+O(1)\) when the shell width is \(\tau\) in scaled
units: the first parameter derivative vanishes at the saddle and the
second is one quarter of the norm-square variance \(O(n/p)\).
The static lower value error adds only \(\epsilon_vk_p\).
Take \(\tau\) and \(\epsilon_v\) sufficiently small.  Both potentials
agree on this shell.  Dividing the preceding upper tail integrals by
this lower bound proves the stated disagreement tails.  Polynomial
moments are absorbed by the same quadratic reference tail and the
subquadratic bound at large radii.
\end{proof}

\begin{lemma}[Positive covariance error]\label{edge:positive-curvature}
For some \(a,c,C>0\), uniformly at every deterministic scale
\(K\le k_p\), \(p\le r_n\), apart from activation points,
\begin{equation}\label{edge:positive-curvature-bound}
 \mathbf P\left(
 \mathcal G\cap
 \{p^2K_p\npreceq(1+Ck_p^{-a})I_{\cH_p}\}\right)
       \le C e^{-c k_p^a}.
\end{equation}
The probability is taken before restricting to \(\mathcal G\).
\end{lemma}

\begin{proof}
Fix a deterministic nonactivation scale \(p\) on the retained spectral
set.  For \(h\in\cH_p\), write \(D_p(h)=p^2K_p(h)-I_{\cH_p}\).
Define a Borel predicate \(\mathcal V_p(W_p,h)\) to be false unless
\(h\in\cH_p\), the saddle \(\alpha_p(h)\) exists, and \(h\ne0\).
Otherwise, with \(e=h/\|h\|\), it requires
\begin{equation}\label{edge:local-parent-validity}
 |\alpha_p(h)|\le s_{\rm aux}p^2,\qquad
 p^2K_p(h)\preceq(1+\epsilon_{\rm aux})I_{\cH_p},\qquad
 p^2\langle e,K_p(h)e\rangle\le0.9.
\end{equation}
This condition contains no restriction to \(\mathcal G\).  At a
nonactivation scale, \(P_p\) is the spectral projection of \(W_p\)
at its flat eigenvalue \(2-p^2\), and
\(R_{p,u}=(2I-W_p+uI)^{-1}\).  Thus the predicate, the saddle, and the
covariance used here are determined by \((p,W_p,h)\).
Their defining formulas are jointly Borel away from activation scales.
In the rest of the proof these objects are evaluated at \(h=h_p\),
and that argument is suppressed.

On \(\mathcal V_p\), use blocks relative to
\(\cH_p=\operatorname{span}\{e\}\oplus(\cH_p\cap e^\perp)\), put
\(z=(D_p)_{\perp e}\), and define on \(\cH_p\cap e^\perp\)
\begin{equation}\label{edge:tangent-error}
 E_p=\left((D_p)_{\perp\perp}
             +\frac{zz^{\mathsf T}}{|(D_p)_{ee}|}\right)_+,
 \qquad \widehat s_p=\|(D_p)_+\|.
\end{equation}
Extend \(E_p\) by zero on the radial direction and on \(\cH_p^\perp\).
Off \(\mathcal V_p\), set both \(E_p\) and \(\widehat s_p\) to zero,
before any restriction to \(\mathcal G\).  Covariance positivity and
\eqref{edge:local-parent-validity} give
\(-I\preceq D_p\preceq\epsilon_{\rm aux}I\) and
\((D_p)_{ee}\le-0.1\) on \(\mathcal V_p\).
Maximizing the quadratic form of \(D_p\) over its radial coordinate,
and then taking the positive part on \(e^\perp\), gives
\begin{equation}\label{edge:tangent-bounds}
 D_p\preceq E_p,\qquad
 \|E_p\|\le C\widehat s_p,\qquad
 \frac{\Tr E_p}{k_p}
       \le C\frac{\|D_p\|_{\rm F}}{\sqrt{k_p}}+\frac C{k_p}
\end{equation}
on \(\mathcal V_p\), with \(D_p\preceq E_p\) understood on \(\cH_p\).
For the norm bound, the maximizing radial coordinate
of a unit transverse vector is bounded because the radial deficit is
fixed and \(D_p\) is bounded; compare the resulting vector with
\(D_p\preceq\widehat s_pI\).  For the trace, the rank-one correction
has bounded trace, and
\(\Tr(((D_p)_{\perp\perp})_+)\le
\sqrt{\dim\cH_p}\|D_p\|_{\rm F}\); the displayed trace bound also
uses \(\dim\cH_p\le Ck_p\) from the retained spectral set.
These are the algebraic bounds of \cite[Lemma~4.4, Equation~(4.4)]{FUA}.
In particular \(\widehat s_p\le\epsilon_{\rm aux}\) and
\(\|E_p\|\le C\epsilon_{\rm aux}\) everywhere, including on disorder failures.
Include \(C\epsilon_{\rm aux}\ll\gamma^2\) in the initial auxiliary choice.

Apply Brascamp--Lieb's covariance inequality \cite{BrascampLieb1976}
to the transition with \(\bar\phi_p\).  After whitening, its Hessian
has a fixed positive margin because
\(\Delta\preceq(1-x^2)p^2I\) and \(\eta\ll\gamma^2\).
If \(\bar K\) is the extended Hessian, the covariance is at most the
transition expectation of
\[
 \Delta^{1/2}
 (I-\Delta^{1/2}\bar K\Delta^{1/2})^{-1}\Delta^{1/2}.
\]
Mollify the semiconcave extension to apply the smooth inequality and
pass to the limit using the positive margin.
Lemma~\ref{edge:annular-extension} replaces the extended Hessian and
transition by \(K_p\) and the true transition on agreement.  Its
polynomially weighted tails control first and second moments on the
complement.  In field units the covariance error is
\(O(p^4M_p^2e^{-c k_p})\); after
\eqref{edge:transition-covariance} and multiplication by \(q^2\),
the remaining polynomial factor in \(k_p\) is absorbed by decreasing
the exponent.

On \(\mathcal G\), the agreement annulus lies in \(\mathcal V_p\), so
\eqref{edge:tangent-bounds} gives
\(K_p\preceq p^{-2}(I+E_p)\).  Put
\(A=I-\Delta/p^2\) and
\(S=p^{-2}\Delta^{1/2}E_p\Delta^{1/2}\).
Then \(A\succeq x^2I\) and
\(\|A^{-1/2}SA^{-1/2}\|\ll1\), so
\[
 (A-S)^{-1}\preceq A^{-1}+2A^{-1}SA^{-1}.
\]
The baseline \(A^{-1}\), substituted in
\eqref{edge:transition-covariance}, is exactly \(q^{-2}I_{\cH_q}\).
For the correction, the two outer factors \(\Delta A^{-1}\) are the
scalar \(\Delta_q/x^2\) on \(\cH_q\).  Thus, after multiplication by
\(q^2\), the first insertion is at most \(2x^{-2}P_qE_pP_q\).
We obtain on \(\mathcal G\), for a child in its annulus,
\begin{equation}\label{edge:parent-inequality}
 \widehat s_q
 \le Cx^{-2}\mathbf E[\|P_qE_pP_q\|\mid W,\cF_q]
       +C_\gamma e^{-c_1k_p}.
\end{equation}
Only \(\Delta_q\) is inverted outside the covariance in
\eqref{edge:transition-covariance}; directions with arbitrarily small
new parent precision cause no loss.

We next compress \(E_p\) before imposing \(\mathcal G\).  The assertion
is trivial when \(E_p=0\); otherwise \(e\) is defined by
\(\mathcal V_p\).  The retained spectral set supplies
\(ck_p\le\dim\cH_p\le Ck_p\) and
\(\dim\cH_q\le Cx^3k_p\); after increasing the threshold,
\(\dim(\cH_p\cap e^\perp)\) is also comparable to \(k_p\).
Condition
under \(\mathbf P\) on the spectrum, \(W_p\) in physical coordinates,
the physical field \(h_p\), and its spectral coordinates.
Equation~\eqref{edge:path-likelihood} makes the remaining rotation in
\(e^\perp\cap\cH_p\) exactly Haar.  The operator \(E_p\) is fixed
under this conditioning: its definition uses only these parent data.
For a uniformly rotating unit vector \(v\) in this space, the Gaussian
representation of the uniform sphere gives
\[
 \mathbf P\left(
 v^{\mathsf T}E_pv>
       \frac{C(\Tr E_p+t\|E_p\|)}{k_p}
       \,\middle|\,\text{parent data}\right)
       \le C(e^{-t}+e^{-c_0k_p}).
\]
Indeed the numerator is a positive Gaussian quadratic form, and the
Gaussian squared norm in the denominator is at least a fixed fraction
of the dimension except with \(e^{-c_0k_p}\) probability.
A fixed-resolution net chosen in the spectral child unit sphere before
the remaining Haar rotation has at most
\(\exp(C\dim\cH_q)\le\exp(Cx^3k_p)\) points.  Project each rotated
net vector onto \(e^\perp\).  Conditional on the parent data, it has
a fixed norm at most one and a uniformly rotating direction.  With
\(t=C_2x^3k_p\), a sufficiently large \(C_2\) pays for the net, and
a sufficiently small \(\gamma\) pays for the second exponential term.
The net estimate for a positive quadratic form gives
\begin{equation}\label{edge:compression}
 \|P_qE_pP_q\|
 \le C\left(\frac{\Tr E_p}{k_p}+x^3\|E_p\|\right)
\end{equation}
outside a set of unconditional \(\mathbf P\)-probability
\(Ce^{-c_\gamma k_p}\).  This is the compression argument of
\cite[Lemma~4.5, Equation~(4.6)]{FUA}; no restriction to \(\mathcal G\)
was made in the Haar law.

Average \eqref{edge:parent-inequality} over
\(I_q=[q/(2\gamma),q/\gamma]\), with normalized logarithmic measure.
On \(\mathcal G\) outside the failure event in \eqref{edge:frobenius-bound},
\eqref{edge:tangent-bounds} and Cauchy--Schwarz give
\[
 \operatorname{Av}_{p\in I_q}\frac{\Tr E_p}{k_p}
       \le C_\gamma k_q^{-b},\qquad b=1/10.
\]
Here \(\operatorname{Av}\) denotes that normalized logarithmic integral.
On its failure the same averaged trace is bounded, because \(E_p\) is
bounded everywhere and \(\dim\cH_p\le Ck_p\).
Its contribution on that failure, with \(\ind_{\mathcal G}\) inserted,
has expectation at most
\(Ce^{-c k_q^{a_1}}\).  The failures of
\eqref{edge:compression} are charged in the same way by the bounded
\(\|E_p\|\).  These estimates are applied to the whole parent-scale
integral before its conditional expectation.  For every nonnegative
error cost \(Y\),
\begin{equation}\label{edge:conditioning-good-event}
 \mathbf E[\ind_{\mathcal G}\mathbf E(Y\mid W,\cF_q)]
       =\mathbf E[\ind_{\mathcal G}Y],
\end{equation}
because \(\mathcal G\) is measurable from \(W\).

The norm term in \eqref{edge:compression} has coefficient
\(Cx^{-2}x^3\le C\gamma\).  The initial choice of \(\gamma\) makes this coefficient
\(\theta<1/2\).  We have proved, on \(\mathcal G\),
\begin{equation}\label{edge:curvature-recurrence}
 \begin{split}
 \widehat s_q&\le
 \theta\operatorname{Av}_{p\in I_q}
       \mathbf E[\widehat s_p\mid W,\cF_q]
       +C_\gamma k_q^{-b}+\operatorname{err}_q,\\
 \operatorname{err}_q&\ge0,\qquad
 \mathbf E[\ind_{\mathcal G}\operatorname{err}_q]
       \le C e^{-c k_q^\beta}
 \end{split}
\end{equation}
for some \(\beta>0\).  Off the child annulus the left side is zero.
The bounded definition of \(E_p\) is essential here: no cost polynomial
in \(n\) is charged on a failed parent interval.

Use a barrier \(B k_q^{-\alpha}\), with \(0<\alpha\le b\) and \(B\)
large.  Since \(k_p=k_qx^{-3}\),
\[
 \theta\operatorname{Av}_{p\in I_q}B k_p^{-\alpha}
       \le\theta(2\gamma)^{3\alpha}B k_q^{-\alpha}.
\]
It absorbs the deterministic error in
\eqref{edge:curvature-recurrence}.  Taking positive parts, applying
conditional Jensen, and using \eqref{edge:conditioning-good-event},
we obtain
\[
 A(q):=\mathbf E[\ind_{\mathcal G}
       (\widehat s_q-Bk_q^{-\alpha})_+]
 \le\theta\operatorname{Av}_{p\in I_q}A(p)+Ce^{-c k_q^\beta}.
\]
For each deterministic nonactivation scale
\(r_n\le p\le C_\gamma r_n\), use
\cite[Proposition~4.2, Equation~(4.2)]{FUA} within its stated range.
Under the biased law \(\mathbf P\), it gives
\(p^2K_p\preceq(1+C(p^{a_0}+k_p^{-a_0}))I\) outside a set whose
intersection with the full event \(\mathcal G_n^{\rm FUA}\) has
probability at most \(Ce^{-n^{\xi_0}}\).
Since \(\mathcal G\subseteq\mathcal G_n^{\rm FUA}\),
the same failure bound holds after intersection with \(\mathcal G\).
Here \(pk_p=np^4\to0\), so \(p^{a_0}\le k_p^{-a_0}\), and
\(k_p\asymp n^{3\cdot10^{-4}}\).  By taking \(\alpha\le a_0\) and a
sufficiently small positive terminal exponential power, the same bound
for \(A(p)\) is \(Ce^{-c k_p^{\beta_0}}\).

Every parent generation multiplies \(k\) by at least
\(\Lambda=(2\gamma)^{-3}>1\).  Iteration, including the terminal
bound, gives for a suitably decreased \(\beta_*>0\)
\[
 A(q)\le C\sum_{j\ge0}\theta^j
       e^{-c(\Lambda^jk_q)^{\beta_*}}
       \le C e^{-c'k_q^{\beta_*}}.
\]
Markov's inequality at the barrier costs only a factor \(k_q^\alpha\),
absorbed by the exponential.  On \(\mathcal G\),
\eqref{edge:local-parent-validity} holds throughout the full auxiliary
annulus, and
\eqref{edge:shrinking-annulus} removes the annulus truncation.
Decrease to a common positive exponent in the barrier and probability.
This proves \eqref{edge:positive-curvature-bound}.
\end{proof}

\begin{proof}[Proof of Proposition~\ref{prop:edge-estimates}]
Lemmas~\ref{edge:spectral} and \ref{edge:norm-density} prove the
spectral and saddle assertions.  Lemmas~\ref{edge:spherical-tails},
\ref{edge:cube-caps}, and \ref{edge:path-annuli} prove the static,
curvature, projection, and ordinary annulus assertions.  Assign a
fixed fraction of the requested failure probability to each of their
spectral and disorder restrictions, and take their finite intersection.
The length bounds give a deterministic \(R_{\rm ann}\) containing the
full annulus; include a further cube-cap request for the static comparisons
on that ball, retaining the previously chosen width \(s_0\).
We select one additional event, independent of \(f\), for the variance.

In the reduction above take \(\eta_p=Ck_p^{-a}\) and let
\(\mathcal N\) be the jointly measurable field set where
\(p^2K_p(h)\npreceq(1+Ck_p^{-a})I_{\cH_p}\).  Set its indicator to one
at activation scales, which have zero scale measure for each
\(W\).  At every deterministic nonactivation scale its evaluation at
\(h_p\) is the failure event in \eqref{edge:positive-curvature-bound}.
Use \(Z_p\) and the common cost \(\rho_p\) from
\eqref{edge:common-failure-cost}; both are independent of \(f\).

Under \(\mathbf P\), conditional Jensen and the spherical identity give
\[
 \mathbf E[\mu_{p,h_p}(Z_p)^2]
       \le\mathbf E[\mu_{p,h_p}(Z_p^2)]
       =\mathbf E[Z_p^2]\le C,
\]
by the uncapped spherical projection tail.  Cauchy--Schwarz and
\eqref{edge:positive-curvature-bound}, absorbing \(k_p\) in the
exponential, imply
\begin{equation}\label{edge:weighted-failure}
 E_U[L_\circ\ind_{\mathcal G}\rho_p]\le C e^{-c k_p^a}.
\end{equation}
The left side is an ordinary path expectation averaged with the Haar
bias.  Take \(q_j=2^j(K/n)^{1/3}\) up to \(r_n\).
Fubini, \eqref{edge:weighted-failure}, and \(L_\circ\ge1/2\) bound
the expected ordinary cost on \([q_j,2q_j]\), intersected with
\(\mathcal G\), by \(Ce^{-c k_{q_j}^a}\).  Markov and
\[
 \sum_{j\ge0}e^{-c(K8^j)^a}\le C e^{-c'K^a}
\]
give an event of arbitrarily large limiting disorder probability on
which, simultaneously in \(j\),
\begin{equation}\label{edge:quenched-error-intervals}
 \int_{q_j}^{\min\{2q_j,r_n\}}\rho_p\,\mathrm d\log p
       \le C e^{-c k_{q_j}^a}.
\end{equation}
Increase \(K\) and use \eqref{edge:haar-bias} for the single discarded
event.  This event is independent of \(f\).

The conditional-variance estimate of
\cite[Proposition~4.1, Equation~(4.1)]{FUA} holds on an ordinary event
of arbitrarily large limiting probability, simultaneously for all real
cube functions and scales in its stated interval \([r_n,p_*]\).  Its
remainder is \(Ce^{-n^{\xi_0}}\|f\|_\infty^2\).  We use it only at
\(r_n\), for \(0\le f\le1\):
\[
 V_{r_n}(f)\le Cr_n^{-2}\cE(f)+Ce^{-n^{\xi_0}}.
\]
Intersect this event with the preceding ones.  In
\eqref{edge:variance-amplification} with terminal scale \(r_n\), its
amplified energy is
\(Cp^{-2}\cE(f)\).  Its amplified remainder is bounded by
\(Ce^{-c k_p^a}\) after decreasing \(a>0\), since
\(k_p\le k_{r_n}=n^{3\cdot10^{-4}}\).
For the other errors use \eqref{edge:quenched-error-intervals}.
The first partial dyadic interval has \(k_{q_j}\asymp k_p\); the rest
increase geometrically.  Hence
\[
 \sum_{q_j\gtrsim p}C(q_j/p)^2e^{-c k_{q_j}^a}
       \le C e^{-c'k_p^a}.
\]
Applying \eqref{edge:variance-amplification} proves
the first inequality of \eqref{edge:variance-estimates}.  It extends
to activation points by continuity of \(V_p\) in \(B_p\).

For the second inequality impose the additional fixed annulus estimate
with \(\zeta_*=\min\{\zeta,s_0/2\}\).  On that annulus the coarse covariance and
\eqref{edge:gradient-bounds} contribute \(CV_p\).  For its complement,
Cauchy--Schwarz and conditional Jensen give
\[
\begin{split}
 k_pQ_\mu[\ind_{\{|\alpha_p(h_p)|>\zeta_* p^2\}}\mu_{p,h_p}(Z_p)]
 &\le k_pQ_\mu\{|\alpha_p(h_p)|>\zeta_* p^2\}^{1/2}
       \mu(Z_p^2)^{1/2}\\
 &\le C e^{-c k_p}.
\end{split}
\]
Here \(\mu(Z_p^2)\le C\) holds at every scale on the chosen uncapped
cube projection event, and the ordinary annulus estimate is also
simultaneous in scale.  Weakening the exponent if needed proves the
second inequality of \eqref{edge:variance-estimates}, on the same
event for every \(f\).
\end{proof}
